%%==========================================================
%%  DOCUMENTO PRINCIPAL
%%  Maestría en Ingeniería de Software
%%  Pontificia Universidad Javeriana
%%  Elaborado por Angela Villota y Luisa Rincón — V 1.0 Abril 2022
%%==========================================================

%%----------------------------------------------------------
%%  PAQUETES
%%----------------------------------------------------------

\documentclass[11pt,twoside,openany]{thesis}
\usepackage{enumerate}
\usepackage{longtable}
\usepackage{url}
\usepackage{hyperref}
\usepackage{makeidx}
\usepackage{graphics}
\usepackage{natbib}
\usepackage[spanish,es-tabla]{babel}
\hbadness=10000

%%----------------------------------------------------------
%%  TÍTULO DEL PROYECTO
%%----------------------------------------------------------

\titulo{Desarrollo de una herramienta CLI para la generación automatizada de proyectos backend e integración de módulos reutilizables}

%%----------------------------------------------------------
%%  AUTORES
%%----------------------------------------------------------

\autorA{codigo}{Jorge Andres Facundo Solarte}

%%----------------------------------------------------------
%%  FECHA Y DIRECTORES
%%----------------------------------------------------------

\fecha{xx de xx de xx}

\directorproyecto{Mg.}{Juan Carlos Martinez Arias}
\directorcarrera{Ph.D.}{Luisa Rincón}

\usepackage{amsmath,amssymb}             % AMS Math
\usepackage{tikz}                        % Diagramas (usado en 1.5 Metodología)
% \usepackage[french]{babel}
%\usepackage[utf8,latin1]{inputenc}
%\usepackage[left=1.0in,right=1.0in,top=1.0in,bottom=1.0in,includefoot,includehead,headheight=13.6pt]{geometry}
\usepackage[left=1.0in,right=1.0in,top=1.0in,bottom=1.0in,includefoot,includehead,headheight=14.33044pt]{geometry}
\raggedbottom                            % Evita que flushbottom estire listas/párrafos para rellenar la página
\renewcommand{\baselinestretch}{1.05}
% Table of contents for each chapter

\usepackage[nottoc, notlof, notlot]{tocbibind}
\usepackage{minitoc}
\setcounter{minitocdepth}{2}
\mtcindent=15pt
% Use \minitoc where to put a table of contents

\usepackage{aecompl}

% My pdf code

\usepackage{ifpdf}


% Links in pdf
\usepackage{color}
\definecolor{linkcol}{rgb}{0,0,0.4} 
\definecolor{citecol}{rgb}{0.5,0,0} 

% Change this to change the informations included in the pdf file

% See hyperref documentation for information on those parameters

\hypersetup
{
bookmarksopen=true,
pdftitle="Registro Calificado Ingeniería de Sistemas y Computación",
pdfauthor="Carrera Ingeniería de Sistemas y Computación", 
pdfsubject="Registro Calificado", %subject of the document
%pdftoolbar=false, % toolbar hidden
pdfmenubar=true, %menubar shown
pdfhighlight=/O, %effect of clicking on a link
colorlinks=true, %couleurs sur les liens hypertextes
pdfpagemode=UseNone, %aucun mode de page
pdfpagelayout=SinglePage, %ouverture en simple page
pdffitwindow=true, %pages ouvertes entierement dans toute la fenetre
linkcolor=linkcol, %couleur des liens hypertextes internes
citecolor=citecol, %couleur des liens pour les citations
urlcolor=linkcol %couleur des liens pour les url
}

% definitions.
% -------------------

\setcounter{secnumdepth}{3}
\setcounter{tocdepth}{2}

% Some useful commands and shortcut for maths:  partial derivative and stuff

\newcommand{\pd}[2]{\frac{\partial #1}{\partial #2}}
\def\abs{\operatorname{abs}}
\def\argmax{\operatornamewithlimits{arg\,max}}
\def\argmin{\operatornamewithlimits{arg\,min}}
\def\diag{\operatorname{Diag}}
\newcommand{\eqRef}[1]{(\ref{#1})}

\usepackage{rotating}                    % Sideways of figures & tables
%\usepackage{bibunits}
%\usepackage[sectionbib]{chapterbib}          % Cross-reference package (Natural BiB)
%\usepackage{natbib}                  % Put References at the end of each chapter
                                         % Do not put 'sectionbib' option here.
                                         % Sectionbib option in 'natbib' will do.
\usepackage{fancyhdr}                    % Fancy Header and Footer

% \usepackage{txfonts}                     % Public Times New Roman text & math font
  
%%% Fancy Header %%%%%%%%%%%%%%%%%%%%%%%%%%%%%%%%%%%%%%%%%%%%%%%%%%%%%%%%%%%%%%%%%%
% Fancy Header Style Options

\pagestyle{fancy}                       % Sets fancy header and footer
\fancyfoot{}                            % Delete current footer settings

%\renewcommand{\chaptermark}[1]{         % Lower Case Chapter marker style
%  \markboth{\chaptername\ \thechapter.\ #1}}{}} %

%\renewcommand{\sectionmark}[1]{         % Lower case Section marker style
%  \markright{\thesection.\ #1}}         %

\fancyhead[LE,RO]{\bfseries\thepage}    % Page number (boldface) in left on even
% pages and right on odd pages
\fancyhead[RE]{\bfseries\nouppercase{\leftmark}}      % Chapter in the right on even pages
\fancyhead[LO]{\bfseries\nouppercase{\rightmark}}     % Section in the left on odd pages

\let\headruleORIG\headrule
\renewcommand{\headrule}{\color{black} \headruleORIG}
\renewcommand{\headrulewidth}{1.0pt}
\usepackage{colortbl}
\arrayrulecolor{black}

\fancypagestyle{plain}{
  \fancyhead{}
  \fancyfoot{}
  \renewcommand{\headrulewidth}{0pt}
}

\usepackage{algorithm}
\usepackage[noend]{algorithmic}

%%% Clear Header %%%%%%%%%%%%%%%%%%%%%%%%%%%%%%%%%%%%%%%%%%%%%%%%%%%%%%%%%%%%%%%%%%
% Clear Header Style on the Last Empty Odd pages
\makeatletter

\def\cleardoublepage{\clearpage\if@twoside \ifodd\c@page\else%
  \hbox{}%
  \thispagestyle{empty}%              % Empty header styles
  \newpage%
  \if@twocolumn\hbox{}\newpage\fi\fi\fi}

\makeatother
 
%%%%%%%%%%%%%%%%%%%%%%%%%%%%%%%%%%%%%%%%%%%%%%%%%%%%%%%%%%%%%%%%%%%%%%%%%%%%%%% 
% Prints your review date and 'Draft Version' (From Josullvn, CS, CMU)
\newcommand{\reviewtimetoday}[2]{\special{!userdict begin
    /bop-hook{gsave 20 710 translate 45 rotate 0.8 setgray
      /Times-Roman findfont 12 scalefont setfont 0 0   moveto (#1) show
      0 -12 moveto (#2) show grestore}def end}}
% You can turn on or off this option.
% \reviewtimetoday{\today}{Draft Version}
%%%%%%%%%%%%%%%%%%%%%%%%%%%%%%%%%%%%%%%%%%%%%%%%%%%%%%%%%%%%%%%%%%%%%%%%%%%%%%% 

\newenvironment{maxime}[1]
{
\vspace*{0cm}
\hfill
\begin{minipage}{0.5\textwidth}%
%\rule[0.5ex]{\textwidth}{0.1mm}\\%
\hrulefill $\:$ {\bf #1}\\
%\vspace*{-0.25cm}
\it 
}%
{%

\hrulefill
\vspace*{0.5cm}%
\end{minipage}
}

\let\minitocORIG\minitoc
\renewcommand{\minitoc}{\minitocORIG \vspace{1.5em}}

\usepackage{multirow}
\usepackage{slashbox}
\usepackage{booktabs}                    % Tablas del Capítulo 3 (estudio comparativo, SRS)
\usepackage{float}                       % Habilita el especificador [H] usado en las tablas nuevas
\usepackage{wasysym}                     % \CIRCLE \LEFTcircle \Circle para la escala de valoración ●/◐/○
\usepackage{graphicx}                    % \rotatebox para encabezados de la tabla comparativa (main.tex solo carga graphics)

% Los captions de longtable usan por defecto \LTcapwidth = 4in, lo que los deja
% en una caja angosta y centrada, distinta a los de las tablas normales.
\setlength{\LTcapwidth}{\textwidth}

\newenvironment{bulletList}%
{ \begin{list}%
	{$\bullet$}%
	{\setlength{\labelwidth}{25pt}%
	 \setlength{\leftmargin}{30pt}%
	 \setlength{\itemsep}{\parsep}}}%
{ \end{list} }

\renewcommand{\epsilon}{\varepsilon}

% centered page environment

\newenvironment{vcenterpage}
{\newpage\vspace*{\fill}\fancyhf{}\renewcommand{\headrulewidth}{0pt}}
{\vspace*{\fill}\par\pagebreak}


\makeindex

\begin{document}

\maketitle

%%----------------------------------------------------------
%%  CARTAS DE PRESENTACIÓN
%%----------------------------------------------------------

\makecoverletter

%%----------------------------------------------------------
%%  FICHA RESUMEN
%%----------------------------------------------------------

\newpage
%%==========================================================
%%  FICHA RESUMEN
%%==========================================================

\thispagestyle{empty}
\begin{center}
    \Large{Ficha Resumen \\ Trabajo de Grado Maestría en Ingeniería de Software}
\end{center}

\textbf{TíTULO:}{Completar}
\begin{enumerate}
    \item Énfasis: Ingeniería de Software
    \item Área de trabajo:
    \item Tipo de proyecto (Aplicado, Innovación, Investigación):
    \item Estudiante:
    \item Correo electrónico:
    \item Dirección y teléfono:
    \item Director:
    \item Vinculación del director: (Planta/Cátedra/Externo)
    \item Correo electrónico del director:
    \item Co-Director (Si aplica):
    \item Grupo o empresa que lo avala (Si aplica):
    \item Otros grupos o empresas:
    \item Palabras clave(al menos 5):
    \item ODS que aplica al proyecto (Agenda 2030): ???
    \item Fecha de inicio:
    \item Resumen:  Debe contener el tema a trabajar en el proyecto de grado, su importancia, la problemática que aborda, los objetivos propuestos, resultados, lecciones aprendidas.  Debe ser redactado en un solo párrafo y no contener espacios entre líneas ni sangría. No debe llevar citaciones.  Debe ser escrito con máximo 400 palabras 
\end{enumerate}


%%----------------------------------------------------------
%%  AGRADECIMIENTOS
%%----------------------------------------------------------

%%==========================================================
%%  AGRADECIMIENTOS
%%==========================================================

\agradecimientos  % No borrar esta línea pues define el formato. OJO el texto debe estar entre llaves
{%NO BORRAR
Escribe aqui los agradecimientos, aquí hay algunas pautas que podrían ser útiles:
\textbf{Comenzar con una nota de agradecimiento general}: Empieza con una declaración general de agradecimiento a todas las personas que te han apoyado a lo largo del proceso.

\textbf{Agradecimientos específicos:}
Familia y amigos: A menudo, las primeras personas que se mencionan son los miembros de la familia y amigos cercanos que han proporcionado apoyo emocional, asesoramientoo.

Asesores y supervisores: Es esencial agradecer a tu asesor de tesis o cualquier otro mentor académico. 
Comité de tesis o jurado: Si aplicable, agradece a los miembros de tu comité de tesis o jurado por su tiempo y atención.

Colegas y compañeros de clase: Menciona a cualquier compañero estudiante o colega que haya contribuido de alguna manera a tu investigación o pensamiento.

Organizaciones de financiación: Si tu investigación fue financiada por alguna beca o subsidio, asegúrate de reconocer la contribución de la organización.

Profesionales y expertos en la materia: Si has tenido ayuda de expertos o profesionales (entrevistas, consultas), menciona su ayuda e influencia en tu trabajo.


Sé sincero y personal, pero mantén la profesionalidad y el enfoque académico de la sección.
Concluye con una declaración que refleje tu experiencia personal a lo largo de tu viaje académico y cómo estos agradecimientos resuenan contigo

} %NO BORRAR


%%----------------------------------------------------------
%%  ABSTRACT / RESUMEN
%%----------------------------------------------------------

%%==========================================================
%%  ABSTRACT / RESUMEN
%%==========================================================

\chapter*{Resumen}
Resumen en español 150 palabras aprox\\
Esta plantilla es una guía para elaborar el trabajo. Puede agregar secciones y hacer los ajustes que requiera para presentar de la mejor manera posible su trabajo.

{\bf Palabras Clave}: entre tres y cinco palabras que diferencien su trabajo. Son las palabras que van a facilitar que tu documento sea visible en cualquier motor de búsqueda. Piense en palabras que describan las ideas principales del
proyecto.

\chapter*{Abstract}
{\bf Keywords}: 

%%----------------------------------------------------------
%%  TABLA DE CONTENIDOS
%%----------------------------------------------------------
  
\newpage{\ } 
\thispagestyle{empty} 
\tableofcontents 
\listoffigures   \listoftables   \mainmatter

\dominitoc
\cleardoublepage

%%----------------------------------------------------------
%%  CAPÍTULO 1: INTRODUCCIÓN
%%----------------------------------------------------------

\chapter{Introducción}
%%==========================================================
%%  CAPÍTULO 1: INTRODUCCIÓN
%%==========================================================

Este capítulo sitúa el trabajo presentando el problema que lo origina, los
objetivos que se comprometió a alcanzar, el alcance dentro del cual se movió,
la metodología que lo organizó y los resultados que produjo.

El trabajo construye una herramienta de línea de comandos que genera
proyectos backend estandarizados y permite incorporarles funcionalidades
transversales reutilizables, para un equipo de desarrollo que opera en una
entidad financiera regulada. Los cuatro capítulos centrales desarrollan cada
una de las fases de la metodología, y los dos últimos cierran el trabajo
contra sus objetivos.

\section{Definición del problema}

\subsection{Planteamiento del problema}
\label{sec:planteamiento}
% Migrado desde Anteproyecto, sección 1 (Planteamiento del Problema)
En los procesos actuales de transformación digital, las organizaciones han aumentado la demanda de desarrollo y evaluación de software con el fin de soportar sus capacidades de negocio. Esta tendencia se enmarca en el reconocimiento de la transformación digital como un impulsor fundamental de la competitividad empresarial, con cerca del 90\% de las organizaciones inmersas en estas iniciativas y un gasto global proyectado que superará los 4 billones de dólares para 2027 \citep{Jie2025}. Este crecimiento ha generado una presión constante por reducir los tiempos de entrega y aumentar la eficiencia de los equipos de desarrollo, debido a que la velocidad de salida al mercado se ha convertido en un factor diferencial para la competitividad y la satisfacción de los usuarios.

Como respuesta a estos desafíos, las metodologías ágiles se han
convertido en un enfoque clave para la gestión y el desarrollo de
proyectos \citep{Palopak2024}. Sin embargo, los equipos enfrentan
tensiones de desempeño como el equilibrio entre velocidad y calidad, lo
que evidencia la persistencia de una brecha entre la velocidad esperada
por el negocio y la capacidad real de mantenerla con calidad técnica
\citep{Armbruster2025}. El Laboratorio Digital del Banco W
\citep{UniversidadIcesi2023} es un caso de esa brecha, pues opera bajo
marcos ágiles con prácticas como Dual Track Agile \citep{MartyCagan2012}
y, a medida que logra entregar más rápido, las expectativas
institucionales aumentan y la agilidad pasa de ser una ventaja a
convertirse en el nuevo estándar exigido.

En su práctica habitual, el inicio de un proyecto implica configurar
desde cero la estructura base del sistema, las convenciones técnicas y
las integraciones con servicios externos, mientras que la incorporación
de funcionalidades transversales entre aplicaciones, como autenticación,
observabilidad o mensajería, se realiza de forma manual reutilizando
código existente. El resultado son funcionalidades duplicadas e
inconsistencias en la estructura de los proyectos que dificultan su
mantenimiento y la evolución controlada del ecosistema de aplicaciones.
Las plantillas base o \emph{boilerplates} buscan encapsular esa
estructura común, pero su reutilización real suele ser limitada porque
tienden a materializarse como artefactos estáticos que quedan
desactualizados o se acoplan a contextos específicos, y se ha señalado
que rara vez logran ser reutilizables de forma general debido a la
diversidad de necesidades y tecnologías entre proyectos
\citep{SeanPromsopeak2022}. Persiste por tanto la necesidad de una base
técnica que permita industrializar el arranque de proyectos y la
composición de capacidades reutilizables sin sacrificar la autonomía de
los equipos ni los atributos de calidad del software.

\subsection{Formulación del problema}
\label{sec:formulacion}

De lo anterior se desprende la pregunta que orienta el trabajo:

\begin{quote}
¿Cómo reducir el esfuerzo de iniciar un proyecto backend y de incorporarle
funcionalidades transversales reutilizables, en un equipo de desarrollo de
una entidad financiera regulada, sin sacrificar la autonomía del equipo sobre
el código que recibe ni los atributos de calidad del software?
\end{quote}

\subsection{Sistematización del problema}
\label{sec:sistematizacion}

La pregunta general se desagrega en cuatro preguntas guía, cada una de las
cuales se responde en un capítulo del trabajo.

\begin{enumerate}
\item ¿Qué necesita realmente un equipo de desarrollo para iniciar un
proyecto backend estandarizado, y qué parte de esa necesidad no resuelve
ninguna de las soluciones disponibles en el mercado?

\item ¿Bajo qué enfoque, qué arquitectura y qué reglas de gobierno debe
estructurarse una familia de proyectos para que un mismo activo pueda servir
a varios de ellos sin que diverjan con el tiempo?

\item ¿De qué manera puede automatizarse la generación de un proyecto
conforme a esa estructura, de modo que el resultado arranque sin
intervención manual y siga perteneciendo al equipo que lo recibe?

\item ¿Cómo puede incorporarse una funcionalidad transversal a un proyecto
ya en curso, y propagarse después un cambio de esa funcionalidad, sin
intervenir proyecto por proyecto ni sobrescribir el código propio de cada
equipo?
\end{enumerate}


\section{Objetivos del proyecto}

\subsection{Objetivo General}
% Migrado desde Anteproyecto, sección 2.1 (Objetivo General)
Desarrollar una herramienta CLI para entornos Node.js (v22–v24) que automatice la generación de proyectos backend bajo principios de arquitectura hexagonal y Domain-Driven Design, alineados con los estándares BIAN e ISO 20022, para reducir el esfuerzo de configuración al iniciar proyectos.

\subsection{Objetivos específicos}
% Migrado desde Anteproyecto, sección 2.2 (Objetivos específicos)
\begin{enumerate}
\item Definir y documentar un marco de referencia que establezca las convenciones, patrones de desarrollo y criterios de organización modular del código, así como los lineamientos de alineación con los estándares BIAN e ISO 20022 que guiarán el desarrollo de la herramienta.

\item Implementar un mecanismo de generación de proyectos mediante plantillas configurables que permita obtener una estructura base ejecutable alineada con el marco de referencia definido.

\item Construir un gestor de módulos reutilizables que permita incorporar funcionalidades transversales mediante contratos de integración definidos y compartidos entre proyectos.

\item Evaluar el comportamiento de la solución desarrollada en el contexto del Laboratorio Digital del Banco W midiendo el tiempo requerido para la generación de un proyecto frente a la línea base de proyectos existentes, la correcta incorporación de un módulo reutilizable y el inicio de implementación de lógica de negocio en el proyecto generado.

\end{enumerate}

\subsection{Divergencias declaradas respecto de los objetivos}
\label{sec:divergencias}

Dos cláusulas de los objetivos transcritos arriba no se cumplieron tal como
están formuladas, y se declaran aquí porque el lector debe saberlo antes de
leer lo que sigue. La primera es el rango de versiones de la plataforma de
ejecución, porque la herramienta exige la versión 24 y rechaza de forma
explícita las dos anteriores, en lugar de soportar el rango completo que el
objetivo enuncia. La segunda es la alineación con ISO 20022, que se revisó
como fundamento del marco de referencia y se excluyó de forma deliberada, con
la razón que la sección~\ref{sec:fundamentos-adoptados} desarrolla.

Ninguna de las dos es un descuido, y las dos tienen una decisión razonada
detrás. El Capítulo~\ref{cap:evaluacion} las recoge, junto con el resto de
las divergencias del proyecto, y da el veredicto de cada objetivo cláusula
por cláusula.

\section{Delimitaciones y alcances}
\label{sec:delimitaciones-alcances}
% Migrado desde Anteproyecto, sección 4 (Alcance del Proyecto)
% Ajuste de tiempo verbal: "se enfocará" -> "se enfocó"; "estará orientado" -> "estuvo orientado"; "se contemplará" -> "se contemplaron"; "quedarán" -> "quedaron"
El proyecto se enfocó en el diseño, implementación y validación de una herramienta CLI orientada a la creación de nuevos proyectos de software a partir de parámetros y configuraciones predefinidas. Además, estuvo orientado principalmente al desarrollo de aplicaciones backend.

Se contemplaron los siguientes componentes y capacidades:

\begin{itemize}
    \item Desarrollo de la herramienta para su uso en un entornos de desarrollo Node.js (v22-v24).
    \item Documentación para su instalación, configuración e implementación.
    \item Generación de una estructura base de proyecto que funcione como plantilla inicial.
    \item Capacidad de iniciar nuevos proyectos a partir de dichas plantillas.
    \item Mecanismo para la implementación de funcionalidades compartidas mediante reglas definidas.
    \item Desarrollo de un sistema de almacenamiento de plantillas, asegurando su disponibilidad para los equipos.
    \item Generación automática de documentación openAPI basada en los módulos y funcionalidades integradas.
\end{itemize}

Quedaron fuera del alcance:

\begin{itemize}
    \item Integración con procesos de despliegue, infraestructura o automatización DevOps.
    \item Implementación de funcionalidades específicas del negocio.
    \item Soporte para otros lenguajes de programación o tecnologías diferentes a Node.js.
\end{itemize}
Más allá de ese alcance funcional, el trabajo operó bajo tres delimitaciones
que condicionan cómo deben leerse sus resultados.

La primera es de contexto. El problema se levantó sobre un único caso de
estudio, el Laboratorio Digital del Banco W, con un equipo técnico de cuatro
desarrolladores. Las necesidades que sostienen la especificación se
contrastaron con desarrolladores de otras organizaciones, pero el contexto de
aplicación es uno solo, y las conclusiones se formulan sobre propiedades
transferibles y no sobre el caso particular.

La segunda es de confidencialidad. El trabajo se desarrolla sobre una
organización del sector financiero, de modo que no se documenta información
privada de la entidad ni del equipo. El contexto se describe por sus
propiedades, las decisiones se sostienen sobre necesidades y requerimientos
verificables, y ningún dato de identificación personal se incorpora al
documento.

La tercera es de disponibilidad del estándar. La alineación con BIAN se
sostiene sobre el documento oficial de cada capacidad de negocio, y el
proyecto dispuso de ese documento para una sola de ellas. Las demás conservan
una descripción derivada de las vistas generales del modelo, y así quedan
declaradas, porque reemplazar una paráfrasis por otra no es transcribir un
estándar.

\section{Metodología de la investigación}
% Migrado desde Anteproyecto, sección 7 (Metodología del proyecto)
% Ajuste de tiempo verbal: "se adopta" -> "se adoptó"; "orientarán/fundamentarán/conformarán" -> "orientaron/fundamentaron/conformaron"; "Diseña y construye" -> "Diseñó y construyó" (Fases 2 y 3); "Evalúa" -> "Evaluó" (Fase 4)
Con el fin de dar cumplimiento al objetivo general propuesto, se adoptó el enfoque de Design Science Research (DSR) como metodología general del proyecto \citep{Wieringa2014}. Este enfoque plantea la construcción de un artefacto que solucione de manera útil y efectiva un problema en un dominio determinado, lo cual se alinea de forma directa con la naturaleza del problema identificado y con los resultados esperados del proyecto.

\begin{figure}[H]
  \scalebox{0.54}{\usetikzlibrary{trees,positioning,arrows.meta,chains,shapes.geometric,shapes,matrix,shapes.symbols}

\definecolor{bgsolid}{RGB}{220,220,220}
\definecolor{bgsub}{RGB}{240,240,240}
\definecolor{border}{RGB}{50,50,50}


\tikzset{
  solid block/.style={
    rectangle, draw=border, thick,
    fill=bgsolid, rounded corners=6pt,
    text width=19em, align=center,
    inner sep=10pt, font=\sffamily\small
  },
  dashed block/.style={
    rectangle, draw=border, thick, dashed,
    rounded corners=6pt,
    align=center,
    inner sep=10pt, font=\sffamily\small
  },
  subnode/.style={
    rectangle, draw=border, solid, thick,
    fill=bgsub, rounded corners=5pt,
    align=center, inner sep=6pt,
    minimum width=2.4cm, font=\sffamily\small
  },
  main arrow/.style={-{Stealth[length=6pt]}, thick, draw=border},
  inner arrow/.style={-{Stealth[length=5pt]}, thick, draw=border},
  inner dashed/.style={-{Stealth[length=5pt]}, thick, draw=border, dashed},
}

\begin{tikzpicture}[node distance=0.7cm]

% FASE 1
\node[solid block] (n1) at (0,0) {
  \textbf{\normalsize Fase 1 — Fundamentación y definición}\\[5pt]
  \textbf{\small Dimensión del problema}\\
  {\footnotesize Caracterización del contexto · Stakeholders}\\
  {\footnotesize Entrevistas · Análisis cualitativo}\\
  {\footnotesize Necesidades y requerimientos}\\[5pt]
  \textbf{\small Dimensión de la solución}\\
  {\footnotesize Antecedentes y soluciones existentes}\\
  {\footnotesize Principios: DDD, Hexagonal}\\
  {\footnotesize Estándares: BIAN, ISO 20022 · Convenciones}
};

% FASE 2
\node[dashed block, below=of n1] (n2) {
  \begin{tikzpicture}[node distance=0.4cm]
    \node[minimum width=16em, minimum height=0pt, inner sep=0pt] (ruler) {};

    \node[font=\sffamily\bfseries\small, align=center,
          below=0pt of ruler] (t2)
      {Fase 2 — Mecanismo de generación de proyectos};

    \node[subnode, below=0.3cm of t2, xshift=-1.5cm] (s2a)
      {\textbf{Diseñar}\\[2pt]
       {\footnotesize Plantilla base}\\
       {\footnotesize Convenciones}};

    \node[subnode, right=0.5cm of s2a] (s2b)
      {\textbf{Prototipar}\\[2pt]
       {\footnotesize Estructura modular}\\
       {\footnotesize Entorno base}};

    \node[subnode, below=0.4cm of s2b] (s2c)
      {\textbf{Refinar}\\[2pt]
       {\footnotesize Ajustar según}\\
       {\footnotesize retroalimentación}};

    \draw[inner arrow, solid] (s2a) -- (s2b);
    \draw[inner arrow] (s2b) -- (s2c);
    \draw[inner dashed] (s2c.west) -| (s2a.south);
  \end{tikzpicture}
};

% FASE 3
\node[dashed block, below=of n2] (n3) {
  \begin{tikzpicture}[node distance=0.4cm]
    \node[minimum width=16em, minimum height=0pt, inner sep=0pt] (ruler) {};

    \node[font=\sffamily\bfseries\small, align=center,
          below=0pt of ruler] (t3)
      {Fase 3 — Gestor de módulos reutilizables};

    \node[subnode, below=0.3cm of t3, xshift=-1.5cm] (s3a)
      {\textbf{Construir}\\[2pt]
       {\footnotesize Gestor de}\\
       {\footnotesize módulos}};

    \node[subnode, right=0.5cm of s3a] (s3b)
      {\textbf{Probar}\\[2pt]
       {\footnotesize Contratos de}\\
       {\footnotesize integración}};

    \node[subnode, below=0.4cm of s3b] (s3c)
      {\textbf{Refinar}\\[2pt]
       {\footnotesize Ajustar según}\\
       {\footnotesize retroalimentación}};

    \draw[inner arrow, solid] (s3a) -- (s3b);
    \draw[inner arrow] (s3b) -- (s3c);
    \draw[inner dashed] (s3c.west) -| (s3a.south);
  \end{tikzpicture}
};

% FASE 4
\node[solid block, below=of n3] (n4) {
  \textbf{\normalsize Fase 4 — Evaluación y validación}\\[4pt]
  Caso de estudio:\\ Laboratorio Digital del Banco W \\ Tiempo de configuración, incorporación de módulo, inicio de implementación
};

\draw[main arrow] (n1) -- (n2);
\draw[main arrow] (n2) -- (n3);
\draw[main arrow] (n3) -- (n4);

\end{tikzpicture}
}
  \caption{Fases de la metodología Design Science Research aplicada al proyecto}
  \label{fig:metodologia}
\end{figure}
\pagebreak

El desarrollo se organiza en cuatro fases (Figura~\ref{fig:metodologia}), las cuales se clasifican en dos tipos según su tipo de ejecución: fases secuenciales y fases iterativas.

\begin{itemize}
  \item \textbf{Fases secuenciales:} Este tipo de ejecución responde a que sus actividades requieren consolidar resultados definidos antes de avanzar, sin necesidad de refinamiento progresivo.

  \item \textbf{Fases iterativas:} Este tipo de ejecución responde a la complejidad técnica inherente a la construcción del artefacto, donde cada iteración permite refinar progresivamente el resultado a partir de los hallazgos obtenidos en la iteración anterior.
\end{itemize}

Esta sección describe la metodología, con el propósito de cada fase y
las actividades que la componen. Su ejecución y los resultados obtenidos
se desarrollan en los capítulos siguientes, dedicando un capítulo a cada
fase: la Fase 1 en el Capítulo~\ref{cap:fundamentacion}, la Fase 2 en el
Capítulo~\ref{cap:generacion}, la Fase 3 en el Capítulo~\ref{cap:gestor}
y la Fase 4 en el Capítulo~\ref{cap:evaluacion}.

\subsubsection*{Fase 1: Fundamentación y definición del marco de referencia}

Estableció las bases conceptuales que orientaron el diseño y construcción del artefacto, mediante el análisis del problema, la revisión de soluciones existentes y la definición de los estándares, convenciones y principios que estructuran la solución.

La fase se abordó en dos dimensiones complementarias. La \textbf{dimensión del problema} (actividades 1 a 4) comprendió un proceso formal de obtención de requisitos sobre un caso de estudio real, con el fin de sustentar en evidencia qué debe resolver la herramienta y bajo qué condiciones sería adoptada, en lugar de asumirlo de antemano. La \textbf{dimensión de la solución} (actividades 5 a 9) comprendió la revisión de las soluciones existentes y la definición del marco de referencia técnico que estructura el artefacto, una vez comprendido el problema.

\paragraph*{Actividades relacionadas}

\begin{enumerate}
    \item Caracterizar el contexto de aplicación en cuanto a su forma de trabajo, su contexto tecnológico y sus prácticas de reutilización de código.

    \item Identificar y justificar los stakeholders del artefacto, y definir el protocolo de entrevistas con guiones diferenciados por rol.

    \item Analizar cualitativamente las entrevistas para identificar los temas recurrentes entre roles y agruparlos en los grupos temáticos que la solución debe atender.

    \item Especificar los requerimientos funcionales y no funcionales de la herramienta, priorizados y trazables a los temas identificados.

    \item Revisar la literatura y antecedentes sobre estandarización y generación de proyectos de software.

    \item Analizar y comparar soluciones existentes en el ámbito de la generación y estandarización de proyectos de software, estableciendo un marco de evaluación e identificando las brechas que ninguna de ellas resuelve.

    \item Identificar y seleccionar los principios y buenas prácticas que fundamentaron la solución, incluyendo arquitectura hexagonal y Domain-Driven Design.

    \item Revisar y adoptar los estándares del sector financiero aplicables, como BIAN e ISO 20022, como referencia para la definición de dominios y capacidades bancarias reutilizables.

    \item Definir las convenciones técnicas y patrones de desarrollo que conformaron la herramienta.
\end{enumerate}

\subsubsection*{Fase 2: Diseño y construcción del mecanismo de generación de proyectos}

Diseñó y construyó de forma iterativa el mecanismo que permite generar proyectos con una estructura modular estandarizada desde el inicio, alineada con el marco de referencia definido en la fase anterior, reduciendo el tiempo de arranque y garantizando coherencia entre proyectos.

\paragraph*{Actividades relacionadas}

\begin{enumerate}
    \item Definir los requisitos y parámetros del mecanismo de generación de proyectos.

    \item Diseñar la estructura modular base alineada con los principios arquitectónicos y estándares bancarios definidos en el marco de referencia.

    \item Construir el mecanismo de generación de proyectos alineado con el marco de referencia definido.

    \item Configurar el entorno de desarrollo de forma reproducible, incluyendo convenciones de código y herramientas base.

    \item Implementar la generación automática de documentación openAPI basada en la estructura del proyecto generado.
\end{enumerate}

\subsubsection*{Fase 3: Diseño y construcción del gestor de modulos reutilizables}

Diseñó y construyó de forma iterativa el sistema que permite incorporar funcionalidades transversales en los proyectos de forma desacoplada, mediante contratos de integración definidos y compartidos entre proyectos, garantizando su reutilización sin generar dependencias entre su implementación concreta y la lógica de negocio de cada proyecto.

\paragraph*{Actividades relacionadas}

\begin{enumerate}
    \item Identificar y categorizar las funcionalidades transversales reutilizables entre proyectos, tomando como referencia los dominios funcionales definidos por BIAN.

    \item Definir los contratos de integración y las reglas de composición que garantizan el desacoplamiento entre módulos y proyectos.

    \item Diseñar e implementar el sistema de composición modular alineado con los contratos de integración definidos y el marco de referencia establecido.

    \item Validar que el sistema de composición modular permite incorporar funcionalidades transversales de forma desacoplada en los proyectos generados.

    \item Verificar que las funcionalidades transversales implementadas cumplen con los lineamientos de alineación con los estándares BIAN e ISO 20022 definidos en el marco de referencia.
\end{enumerate}

\subsubsection*{Fase 4: Evaluación y validación del artefacto}

Evaluó la solución desarrollada mediante un caso de estudio en el Laboratorio Digital del Banco W, midiendo el tiempo de configuración inicial frente a proyectos existentes como línea base, la correcta incorporación de un módulo reutilizable y el inicio de implementación de lógica de negocio en el proyecto generado.

\paragraph*{Actividades relacionadas}

\begin{enumerate}
    \item Definir los criterios de evaluación del caso de estudio, incluyendo la medición del tiempo de configuración inicial frente a la línea base de proyectos existentes, la incorporación de un módulo reutilizable y el inicio de implementación de lógica de negocio en el proyecto generado.

    \item Aplicar la solución en un caso de estudio en el Laboratorio Digital del Banco W, generando un proyecto base e incorporando un módulo reutilizable en un entorno de desarrollo real.

    \item Recolectar y analizar los datos obtenidos sobre el tiempo de configuración inicial, la incorporación del módulo reutilizable y el inicio de implementación de lógica de negocio, comparando los resultados frente a la línea base establecida.

    \item Elaborar el informe de evaluación con los hallazgos y recomendaciones sobre la pertinencia del artefacto.
\end{enumerate}

\section{Resultados obtenidos}

El trabajo produjo cuatro resultados, uno por fase, y cada uno se desarrolla
en el capítulo correspondiente.

\begin{enumerate}
\item \textbf{Un marco de referencia técnico} que encuadra la solución como
línea de producto de software y fija la arquitectura de la familia de
proyectos que la herramienta genera, con sus fundamentos adoptados, la
organización del código, los contratos de integración y el modelo de
gobierno de los activos. Es el insumo del que depende todo lo demás, y se
desarrolla en el Capítulo~\ref{cap:fundamentacion}.

\item \textbf{Un mecanismo de generación de proyectos}, materializado en una
herramienta de línea de comandos instalable que produce un proyecto backend
ejecutable sin intervención manual, con las convenciones del marco aplicadas,
el entorno de trabajo configurado y su documentación de API generada bajo el
estándar OpenAPI. Junto a él, un almacén de plantillas en el que una
plantilla es un activo que se describe a sí mismo, con dos plantillas
publicadas que derivan formas de producto distintas. Se desarrolla en el
Capítulo~\ref{cap:generacion}.

\item \textbf{Un gestor de módulos reutilizables} que incorpora una
funcionalidad transversal a un proyecto ya en curso, informa de su estado, la
actualiza preservando el código propio del equipo y permite retirarla, todo a
través de contratos de integración declarados en un artefacto neutral. El
catálogo cuenta con dos módulos construidos, de naturaleza distinta entre sí.
Se desarrolla en el Capítulo~\ref{cap:gestor}.

\item \textbf{Un informe de evaluación} del artefacto construido, medido con
las personas del contexto de aplicación y contra una línea base establecida
antes de que la herramienta existiera. Se desarrolla en el
Capítulo~\ref{cap:evaluacion}.
\end{enumerate}

Los activos de la línea quedaron publicados en el registro de paquetes de la
organización, que es la condición sin la cual ninguno de los tres primeros
resultados sería utilizable por un equipo distinto del que los construyó.

%%==========================================================
%%  RIESGOS ÉTICOS
%%  Se incluye como sección final del Capítulo 1.
%%==========================================================

\section{Riesgos éticos}
\label{sec:riesgos-eticos}

Un sistema que estandariza y automatiza el inicio de aplicaciones concentra
en un único punto decisiones técnicas que después se replican en varios
proyectos. En un contexto financiero esa concentración importa más, por el
manejo de información sensible y de lógica de negocio regulada. El proyecto
identificó cuatro riesgos éticos al plantearse, cada uno con la mitigación a
la que se comprometía, y lo que corresponde decir aquí es cuál de las cuatro
se cumplió. El veredicto final de cada uno, con la evidencia de la
evaluación, lo da el Capítulo~\ref{cap:evaluacion}.

\paragraph*{Exposición de información sensible en plantillas y módulos}

El compromiso era separar de forma estricta la estructura base de los valores
de configuración sensibles, promoviendo variables de entorno y mecanismos de
gestión de secretos desde el diseño. \textbf{Se cumplió, y quedó verificado
por ejecución.} Es lo que la especificación de requerimientos exige como
atributo de configurabilidad y lo que el artefacto hace, porque ninguna
dirección ni credencial vive en el código del producto, toda
configuración entra por variable de entorno validada contra un esquema, el
circuito de verificación barre el código del producto derivado buscando
valores fijos y direcciones, y el registro de eventos del producto redacta
por nombre de clave todo secreto o identificador de persona. La verificación
automática es más de lo que el compromiso pedía.

\paragraph*{Propagación de lógica financiera mal interpretada}

El compromiso era someter a revisión técnica especializada los módulos con
lógica financiera antes de incorporarlos al catálogo. \textbf{El riesgo no se
materializó}, porque ninguno de los dos módulos construidos lleva lógica
financiera, ya que uno guarda y lee registros y el otro exporta señales de
observabilidad. Las dos capacidades son técnicas y transversales por decisión
del marco, que reparte lo técnico al módulo y el dominio al producto. Lo que
queda escrito es la condición de reactivación, porque el día que el catálogo
reciba un módulo que interprete un estándar financiero, y el de mensajería
excluido en la sección~\ref{sec:fundamentos-adoptados} es el caso obvio, la
revisión especializada es exigible antes de su entrada.

\paragraph*{Exclusión de equipos con menor capacidad técnica}

El compromiso tenía dos partes y \textbf{se cumplió a medias}. La primera era
elaborar documentación accesible orientada a distintos niveles de
experiencia, y \textbf{no se construyó}. La propuesta existió, y se descartó
con un argumento de oportunidad, que esa documentación crece con los módulos
incorporados y hacerla antes documenta un producto que todavía va a cambiar
de forma. El argumento es razonable como decisión de alcance de una fase y no
es una respuesta al riesgo ético, porque nadie advirtió que lo fuera, ya que
el rechazo comparó la propuesta contra el calendario y no contra el
compromiso. Se declara así, y la documentación por niveles queda como
recomendación de adopción.

La segunda parte sí está. El producto derivado nace con documentación que
explica qué es, qué hace cada operación, cuál lleva efecto de negocio y cuál
es un andamio; el recorrido de creación se lee en lenguaje de negocio con el
nombre del estándar al lado; la herramienta declara su alcance antes de
escribir y pide confirmación; y los mensajes de error dicen qué encontró, qué
esperaba y qué resuelve la situación. Todo eso baja la barrera de adopción,
que es lo que el riesgo persigue. El acompañamiento con procesos de formación
y transferencia de conocimiento es organizacional y queda fuera de lo que un
artefacto puede resolver.

\paragraph*{Concentración de poder técnico en quienes controlan el sistema}

El compromiso era definir un modelo de gobernanza distribuido con procesos
participativos, conocimiento documentado de forma exhaustiva y una evolución
que responda a necesidades colectivas. \textbf{El mecanismo existe y el
ejercicio no}, y la razón es una condición del entorno declarada desde la
segunda fase.

El mecanismo está completo, porque el catálogo declara un responsable por
activo, la aprobación del gobierno de la línea es obligatoria sobre la ruta
del manifiesto, y el procedimiento de incorporación tiene sus pasos y sus
condiciones de validez verificables por tipo de activo. La documentación
exhaustiva es la parte más fuerte del compromiso, con casi doscientas
decisiones registradas con su contexto, su alternativa y su traza, más de
treinta documentos de expediente y un marco con su propio mecanismo de
evolución, que es lo que permite que otra persona retome el sistema. Y la
evolución respondió a necesidades que no son del autor, porque las doce
necesidades y las tres brechas salen de las entrevistas con las personas del
contexto, y varias rondas de mejora las abrieron una revisión externa y
reportes de uso.

Lo que no hay es más de un participante. El proyecto se construyó con una
sola persona en los tres roles, de modo que proponente y revisor coinciden, y
una gobernanza distribuida instanciada con un participante es un mecanismo
probado y un ejercicio ausente. Trasladarlo a operación consiste en cambiar
quién figura como responsable de cada activo, lo que el modelo de gobierno ya
contempla, y queda como recomendación de adopción.


\section{Resumen del capítulo}

Este capítulo estableció el problema, los objetivos y el marco dentro del
cual se movió el trabajo. El problema es que iniciar un proyecto backend y
dotarlo de funcionalidades transversales es, en el contexto estudiado, un
proceso manual que se concentra en pocas personas y que produce proyectos
que divergen entre sí. La respuesta comprometida es una herramienta que
genere proyectos estandarizados y permita componerles capacidades
reutilizables, construida bajo Design Science Research en cuatro fases.

Dos cosas quedaron declaradas aquí, y las dos deben retenerse al leer lo que
sigue. La primera son las dos divergencias respecto de los objetivos
comprometidos, el piso de la plataforma de ejecución y la exclusión de uno de
los dos estándares financieros, que se nombran antes de desarrollarlas en
lugar de dejarlas para el cierre. La segunda son las delimitaciones bajo las
que el trabajo operó, en particular que el problema se levantó sobre un único
caso de estudio y que la alineación con el estándar de negocio se sostiene
sobre el documento oficial de una sola capacidad.

Los cuatro capítulos siguientes desarrollan una fase cada uno, en el orden en
que se ejecutaron, y el séptimo cierra el trabajo contra los objetivos que
este capítulo transcribe. El Capítulo~\ref{cap:marco} fija antes el
vocabulario y revisa lo que existe.


%%----------------------------------------------------------
%%  CAPÍTULO 2: MARCO TEÓRICO Y ESTADO DEL ARTE
%%----------------------------------------------------------

\chapter{Marco teórico y estado del arte}
%%==========================================================
%%  CAPÍTULO 2: MARCO TEÓRICO Y ESTADO DEL ARTE
%%==========================================================

\label{cap:marco}

\section{Marco teórico}
% Migrado desde Anteproyecto, sección 5.1 (Bases Teóricas)
La revisión conceptual se organiza en cuatro grupos que avanzan de lo
general a lo específico, desde los fundamentos de la disciplina hasta los
estándares propios del dominio financiero en el que opera el contexto de
aplicación.

\subsection{Fundamentos de la ingeniería de software}

Las primeras entradas fijan el vocabulario general sobre el que se apoya
el resto del capítulo, desde la disciplina misma hasta las piezas de
software reutilizable con las que se construye un producto.
\paragraph*{Ingeniería de Software}
Enfoque sistemático y disciplinado para el diseño, desarrollo y mantenimiento de sistemas de software. A través de la aplicación de principios de ingeniería y buenas prácticas, esta disciplina busca construir soluciones confiables, eficientes y escalables que respondan a necesidades específicas de los usuarios o de una organización. Más allá de la implementación de código, la ingeniería de software abarca diversas actividades a lo largo del ciclo de vida de un sistema, incluyendo el análisis de requerimientos, el diseño de la arquitectura, la implementación, las pruebas y el mantenimiento continuo del software, con el fin de garantizar su correcto funcionamiento y evolución en el tiempo \citep{Jazayeri2003}.
\paragraph*{Ciclo de vida del Software}
Proceso estructurado utilizado para planificar, desarrollar, probar, desplegar y mantener sistemas de software. Este enfoque organiza el desarrollo en una serie de fases que abarcan desde la conceptualización inicial de un producto hasta su implementación, operación, mantenimiento y eventual retiro. A través de esta estructura, el SDLC permite gestionar de manera ordenada las actividades involucradas en la construcción de software, asegurando que cada etapa produzca los resultados necesarios para la siguiente. De esta forma, el modelo facilita el control del proceso de desarrollo, mejora la eficiencia del trabajo de los equipos y contribuye a garantizar la calidad del software entregado \citep{Jazayeri2003}.
\paragraph*{Framework de desarrollo}
Conjunto estructurado de componentes de software reutilizables que provee una base y un conjunto de reglas para crear aplicaciones de forma más eficiente, consistente y estandarizada, facilitando el desarrollo mediante módulos predefinidos, APIs, librerías y herramientas que aceleran la construcción de funcionalidades sin necesidad de programar desde cero elementos comunes, como manejo de datos, control de flujo, seguridad o interfaces de usuario. Además, un framework ayuda a mantener la calidad, coherencia y consistencia del código, promueve buenas prácticas de desarrollo y facilita la organización y colaboración dentro de los equipos de desarrollo a lo largo del ciclo de vida de un proyecto \citep{aws-framework-2026}.

\subsection{Diseño y arquitectura de software}

El segundo grupo reúne los principios de diseño que determinan cómo se
organiza internamente un proyecto y qué separa su lógica de negocio de
la tecnología que la soporta. Son los que el marco de referencia de la
solución adopta más adelante como criterios de estructura.
\paragraph*{Modularidad en sistemas de Software}
Principio de desarrollar un sistema mediante la división en subsistemas más pequeños que pueden diseñarse e implementarse de manera relativamente independiente. Estos subsistemas pueden adoptar diferentes formas según el paradigma de desarrollo utilizado, como funciones en sistemas estructurados, clases en sistemas orientados a objetos, aspectos en programación orientada a aspectos o componentes en arquitecturas basadas en componentes. En términos generales, todos estos elementos pueden entenderse como unidades funcionales que encapsulan una parte específica de la lógica del sistema \citep{Yu2009}.
\paragraph*{Domain Driven Design}
Enfoque que propone centraliza el desarrollo de software en el dominio del negocio, buscando que los modelos reflejen con precisión los procesos y objetivos críticos de la organización \citep{Vernon2013}. Este enfoque facilita la colaboración entre expertos de negocio y desarrolladores mediante un lenguaje común y promueve la estructuración de la aplicación en módulos y contextos delimitados con responsabilidades claras, lo que permite modularidad, desacoplamiento y reutilización de funcionalidades. Además, DDD orienta el diseño de la arquitectura y sus componentes para que estén alineados con el dominio, mejorando la mantenibilidad, la coherencia técnica y la capacidad de evolución controlada del sistema.
\paragraph*{Arquitectura hexagonal}
Enfoque arquitectónico que propone organizar las aplicaciones de manera que la lógica central del sistema permanezca independiente de las tecnologías externas.

En este modelo, el núcleo de la aplicación contiene el dominio del sistema, es decir, las reglas de negocio y comportamientos que representan la funcionalidad principal de la aplicación. Al mantener este núcleo aislado de elementos como interfaces de usuario, bases de datos o servicios externos, se busca evitar que las decisiones tecnológicas afecten directamente la lógica del negocio \citep{AlistairCockburn2005}.

Para lograr esta separación, la arquitectura define puertos, que representan interfaces mediante las cuales el núcleo del sistema expone sus funcionalidades, y adaptadores, que son componentes encargados de conectar estas interfaces con tecnologías específicas. De esta manera, los adaptadores traducen las interacciones provenientes del exterior hacia las operaciones definidas por el dominio de la aplicación, permitiendo desacoplar la lógica de negocio de los detalles de infraestructura y facilitando la evolución y mantenimiento del sistema \citep{AlistairCockburn2005}.

\subsection{Líneas de producto de software}

El tercer grupo corresponde al enfoque de ingeniería bajo el cual se
encuadra la solución propuesta. Define qué es una familia de productos
construida sobre activos comunes, cómo se separa la construcción de esos
activos de la de cada producto, cómo se representa lo que varía entre
ellos, y qué ocurre cuando esa reutilización se resuelve copiando.
\paragraph*{Línea de Producto de Software}
Conjunto de sistemas intensivos en software que comparten un conjunto común y gestionado de características, que satisfacen las necesidades de un segmento determinado del mercado o de una misión, y que se desarrollan a partir de un conjunto común de activos siguiendo un procedimiento establecido \citep{Clements2001}. El enfoque surge de reconocer que las similitudes entre los productos de una organización pueden tratarse como una ventaja de ingeniería, y no como una coincidencia que cada proyecto resuelve por su cuenta.

De la definición se desprenden cuatro elementos constitutivos, cuya ausencia impide hablar de una línea de producto. El primero es un alcance gestionado, es decir, la declaración explícita de qué productos pertenecen a la familia; sin él los activos se generalizan hasta volverse inservibles o se especializan hasta servir a un solo producto. El segundo es una arquitectura común compartida por todos los productos, que es lo que hace posible que un mismo activo sirva a varios. El tercero son los activos centrales o \emph{core assets}, que incluyen arquitectura, contratos, modelos y documentación, y no solo código. El cuarto es un mecanismo de derivación, entendido como el procedimiento que produce un producto concreto a partir de los activos y de una selección de características \citep{Clements2001}.

Una línea de producto no es una categoría de herramienta sino un
enfoque de ingeniería, y esa distinción es la que evita el error más
frecuente al interpretar la literatura. Los generadores de
plantillas, los gestores de paquetes y las plataformas internas de
desarrollo son mecanismos de reutilización. Cada uno atiende una
preocupación concreta dentro de su ámbito, pero ninguno declara
alcance, arquitectura ni procedimiento, porque no es su propósito.
Una línea de producto no los sustituye ni compite con ellos, sino
que opera en un plano distinto, los ordena y se materializa
empleándolos \citep{Clements2001}.
\paragraph*{Ingeniería de dominio e ingeniería de aplicación}
Construir los activos reutilizables y construir un producto con ellos son actividades distintas, con insumos, salidas y ritmos diferentes, por lo que la disciplina las separa en dos ciclos. La ingeniería de dominio toma el conocimiento del dominio como insumo y produce los activos, incluidos el modelo de variabilidad y los componentes reutilizables. La ingeniería de aplicación toma los requerimientos de un producto concreto, consulta esos activos y produce el producto \citep{Pohl2005}.

La relación entre ambos ciclos no es de un solo sentido. La ingeniería de aplicación consume los activos, pero también puede revelar una necesidad que el dominio no había contemplado y que debe incorporarse a los activos comunes, quedando así disponible para todos los productos futuros \citep{Pohl2005}.
\paragraph*{Variabilidad y modelo de características}
La variabilidad es la propiedad de un activo de admitir más de una realización, y un punto de variación es el lugar concreto donde esa elección ocurre. El mecanismo formal para representar qué varía entre los productos de una familia y qué permanece fijo es el modelo de características o \emph{feature model}, introducido por el método de análisis de dominio orientado a características \citep{Kang1990}. Un modelo de este tipo organiza las características en forma de árbol y declara, para cada nodo, si es obligatoria, opcional, alternativa o acumulativa respecto de su padre, junto con las restricciones que cruzan el árbol \citep{Apel2013}.

Una característica no solo se selecciona, sino que también se
resuelve en algún momento del ciclo de vida del producto, y ese
momento se denomina tiempo de ligado o \emph{binding time}. Una
característica resuelta durante la generación queda escrita en el
código y cambiarla exige volver a derivar el producto; una resuelta
en tiempo de ejecución puede modificarse mediante configuración
externa sin regenerar nada. El ligado temprano produce productos más
simples y verificables a costa de flexibilidad posterior, mientras
que el ligado tardío conserva flexibilidad a costa de exigir
mecanismos de configuración y validación \citep{Apel2013}.
\paragraph*{Clone-and-own}
Construcción de un producto nuevo a partir de la copia de uno existente, sin que se conserve vínculo alguno entre la copia y su origen. La literatura lo reconoce como una forma legítima de reutilización en etapas tempranas, por su bajo costo inicial y su nula fricción de adopción \citep{Dubinsky2013}.

El problema que la literatura documenta de manera característica es
la divergencia progresiva. Al no existir un mecanismo que propague
los cambios, cada copia evoluciona por su cuenta, lo que incrementa
el esfuerzo de propagación, de análisis de dominio y de control de
calidad, y genera conflictos cuando el origen y las copias
evolucionan en paralelo. El costo de esta práctica crece con el
número de copias \citep{Dubinsky2013}.

\subsection{Estándares del sector financiero}

El último grupo recoge los estándares del dominio en el que opera el
contexto de aplicación, que delimitan cómo se nombran y se organizan las
capacidades de negocio de una entidad bancaria y cómo se estructura el
intercambio de información entre instituciones.
\paragraph*{BIAN (Banking Industry Architecture Network)}
Asociación global independiente y sin fines de lucro formada por instituciones financieras, proveedores de software, consultoras, integradores y académicos con el propósito de definir y promover un marco arquitectónico común para la banca. La misión de BIAN es ofrecer un modelo estandarizado de referencia para la interoperabilidad de sistemas bancarios, apoyando la innovación tecnológica, reduciendo los costos de integración y simplificando la arquitectura de TI en el sector financiero. Para lograr esto, BIAN reúne experiencia de su comunidad para crear un marco de mejores prácticas que los bancos pueden utilizar al diseñar, implementar y probar servicios y plataformas financieras \citep{bian-2026}.
\paragraph*{ISO 20022}
Estándar internacional de mensajería financiera desarrollado por el Technical Committee ISO/TC68 Financial Services de la Organización Internacional de Normalización (ISO). Este estándar proporciona una plataforma común para el desarrollo de mensajes financieros mediante un conjunto de reglas, un diccionario de datos y un catálogo de modelos que describen los procesos de negocio y los elementos de datos utilizados en las comunicaciones financieras. Los modelos y los mensajes derivados se publican en el ISO 20022 Financial Repository, lo que permite a las instituciones financieras definir conjuntos de mensajes basados en un enfoque internacionalmente acordado y consistente en cuanto a semántica y estructura \citep{iso20022-2026}.

\section{Estado del arte}
\label{sec:estado-del-arte}
% Migrado desde Anteproyecto, sección 5.2 (Antecedentes)
En el desarrollo de software, se han implementado diversas estrategias y herramientas para agilizar el inicio de proyectos y garantizar consistencia en las prácticas de desarrollo. Estas soluciones pueden organizarse jerárquicamente según su alcance y complejidad, desde los fundamentos conceptuales hasta plataformas completas de gestión del ciclo de vida de aplicaciones.

\subsection{Fundamentos conceptuales}
\paragraph*{Scaffolding}
Técnica que permite generar automáticamente la estructura básica de una aplicación o de determinados componentes dentro de un proyecto. A partir de configuraciones o modelos definidos, las herramientas de scaffolding crean archivos, directorios y código inicial necesarios para implementar funcionalidades comunes. Esto permite a los desarrolladores contar con una base funcional sobre la cual continuar el desarrollo, acelerando las primeras etapas de construcción del sistema.

Diversos frameworks modernos incorporan herramientas de scaffolding que permiten generar automáticamente módulos o componentes predefinidos, reduciendo el tiempo necesario para implementar nuevas funcionalidades. Por ejemplo, el framework NestJS proporciona una herramienta de línea de comandos (CLI) que permite generar elementos como módulos, controladores y servicios mediante comandos específicos, simplificando la creación de estructuras estándar dentro de las aplicaciones \citep{nestjs-2026}.
\paragraph*{Boilerplate}
Fragmentos de código que pueden reutilizarse en diferentes contextos con pocas o ninguna modificación. Este término proviene originalmente del uso de plantillas en documentos, donde se reutilizan textos estándar que luego se adaptan ligeramente a situaciones específicas. De forma similar, en el desarrollo de software los programadores utilizan secciones de código previamente escritas para evitar tener que implementarlas nuevamente desde cero \citep{aws-boilerplate-2026}.

El uso de boilerplate permite mejorar la eficiencia del desarrollo, ya que facilita la reutilización de código probado y reduce el tiempo necesario para implementar funcionalidades comunes. Además, contribuye a mantener consistencia en la estructura del software y disminuir errores derivados de la repetición manual de código \citep{aws-boilerplate-2026}.

\subsection{Herramientas especializadas de generación de código}
\paragraph*{Yeoman}
Herramienta de código abierto diseñada para facilitar la creación de nuevos proyectos mediante el uso de generadores. Estos generadores actúan como complementos que pueden ejecutarse desde la línea de comandos para generar automáticamente proyectos completos o partes específicas de una aplicación.

El objetivo de Yeoman es ayudar a los desarrolladores a iniciar proyectos rápidamente, proporcionando estructuras iniciales, configuraciones y herramientas que incorporan buenas prácticas de desarrollo \citep{yeoman-2026}.
\paragraph*{Cookiecutter}
Herramienta de generación de proyectos a partir de plantillas, agnóstica de lenguaje, que permite a un equipo partir siempre de una base consistente y reducir la configuración manual repetida. 
    
Sus plantillas son directorios con una raíz parametrizada y un archivo de declaración de variables. Al ejecutarse, la herramienta sustituye esas variables tanto en el contenido como en los nombres de archivos y carpetas mediante el motor Jinja2, y admite condicionales, bucles y scripts de validación antes y después de la generación \citep{cookiecutter-2023}.
\paragraph*{Cruft}
Complemento construido sobre Cookiecutter que resuelve la desalineación entre una plantilla y los proyectos generados a partir de ella.
    
Para lograrlo, registra en el proyecto la versión exacta de la plantilla utilizada, lo que permite comprobar en cualquier momento si el proyecto quedó desactualizado. Al solicitar una actualización, la herramienta aplica los cambios pendientes mediante un proceso de fusión similar al de un sistema de control de versiones \citep{cruft-2026}.
\paragraph*{Copier}
Generador de proyectos a partir de plantillas que mantiene un vínculo persistente entre la plantilla y los proyectos ya generados. Cada proyecto registra la versión de plantilla con la que fue creado, y su comando de actualización reconstruye el proyecto tal como era con esa versión, calcula las modificaciones que introdujo el desarrollador, genera el resultado con la versión nueva de la plantilla y reaplica encima esos cambios, señalando los conflictos que no puede resolver de forma automática \citep{copier-2026}.
\paragraph*{CLI de NestJS}
Herramienta de línea de comandos incorporada al framework NestJS, construida sobre un sistema de \emph{schematics} que describe cómo debe generarse cada elemento. Además de crear los archivos del bloque solicitado, ya sea un módulo, un controlador, un servicio o un recurso completo, modifica el código existente para integrarlo, por ejemplo registrando un servicio recién generado en el arreglo de \emph{providers} de su módulo \citep{nestjs-2026}.
\paragraph*{Nx.js}
Plataforma de desarrollo de código abierto orientada a la gestión de repositorios monolíticos (monorepos) y a la automatización de tareas dentro de proyectos de software. Esta herramienta proporciona un sistema de construcción que permite ejecutar y coordinar tareas como compilación, pruebas y análisis de código de manera eficiente, especialmente en proyectos que contienen múltiples aplicaciones o librerías dentro de un mismo repositorio \citep{nx-2026}.

Además, la herramienta puede generar aplicaciones, librerías y componentes mediante plugins y generadores, facilitando la estandarización de la arquitectura y la organización del código en proyectos de gran escala \citep{nx-2026}.

\subsection{Plataformas de desarrollo internas (IDP)}
\paragraph*{Backstage}
Framework de código abierto para organizar las herramientas, servicios y documentación utilizados en proyectos de software mediante portales de desarrollo. Originalmente creado por Spotify para enfrentar la creciente complejidad de su infraestructura interna, Backstage permite a los equipos desarrollar y gestionar software de manera más eficiente al proporcionar un punto único de interacción con todos los recursos técnicos de la organización \citep{backstage-2026}.

Su base es un catálogo de software centralizado que contiene metadatos sobre servicios, librerías, pipelines de datos y otros componentes, facilitando su descubrimiento y administración. Además, incluye plantillas de software que permiten generar nuevos proyectos de forma estandarizada según las mejores prácticas de la organización, así como un sistema de documentación técnica integrado (TechDocs) que sigue un enfoque de "docs like code" \citep{backstage-2026}.
\paragraph*{Port}
Plataforma diseñada para unificar y organizar todas las herramientas, datos y flujos de trabajo que los equipos de ingeniería utilizan durante el ciclo de vida del software. Es una solución enfocada en mejorar la experiencia del desarrollador proporcionando un punto centralizado donde encontrar, lanzar, operar y automatizar tareas relacionadas con el desarrollo y la infraestructura de software \citep{port-platform-2026}.

Integra varios componentes clave que ayudan a gestionar el ecosistema de desarrollo de forma coherente y eficiente. Entre ellos se encuentran un catálogo de software que unifica metadatos de servicios, entornos y otros activos técnicos, acciones preconfiguradas que estandarizan operaciones repetitivas, un orquestador de flujos de trabajo que automatiza procesos y scorecards para medir la calidad, seguridad y cumplimiento de los estándares del equipo \citep{port-platform-2026}.

\section{Resumen del capítulo}

Este capítulo fijó el vocabulario y revisó lo que existe. El marco teórico
avanzó de lo general a lo específico, desde los fundamentos de la disciplina
hasta los dos cuerpos conceptuales de los que el trabajo se sirve de forma
directa, que son las líneas de producto de software, como enfoque que trata
la reutilización planificada entre productos de una misma familia, y los
estándares del sector financiero, que aportan el vocabulario de negocio con
que se delimita y se nombra un producto.

El estado del arte describió ocho soluciones, organizadas por alcance. Las
herramientas de plantilla producen el punto de partida de un proyecto nuevo;
las que sostienen su ciclo de vida registran de qué versión de plantilla
proviene y propagan los cambios posteriores; y las plataformas de desarrollo
interno catalogan los activos de una organización y les asignan propiedad.
Cada una resuelve bien su propia preocupación.

Lo que este capítulo deliberadamente no hace es evaluarlas. Describir una
herramienta y medirla contra un problema son dos actos distintos, y el
segundo exige antes haber establecido cuál es el problema. Esa evaluación, y
las tres brechas que de ella resultan, corresponden al
Capítulo~\ref{cap:fundamentacion}, que las mide contra los temas
levantadas en el contexto de aplicación.

%%----------------------------------------------------------
%%  CAPÍTULO 3: ANÁLISIS DEL PROBLEMA Y MARCO DE REFERENCIA DE LA SOLUCIÓN
%%----------------------------------------------------------

\chapter{Análisis del problema y marco de referencia de la solución}
%%==========================================================
%% CAPÍTULO 3: ANÁLISIS DEL PROBLEMA Y MARCO DE REFERENCIA DE LA SOLUCIÓN
%%==========================================================
\label{cap:fundamentacion}

La primera fase de la metodología del proyecto (Figura~\ref{fig:metodologia})
persiguió dos propósitos, el primero de los cuales fue demostrar con evidencia
que el problema existe y establecer en qué consiste exactamente, y el segundo
definir cómo debía estructurarse la solución antes de escribir la primera
línea de código. Al cerrarla quedó establecido qué debe hacer la herramienta y
qué forma deben tener los proyectos que produzca.

El trabajo se organizó en dos bloques, y el primero reunió la
evidencia que sustenta el proyecto, levantando el contexto donde la
herramienta iba a operar, el testimonio de quienes trabajan en él y lo que
ofrecen las soluciones existentes, hasta especificar formalmente qué debía
construirse. El segundo definió el encuadre de la solución y los fundamentos
sobre los que se apoya, que es diseño de alto nivel y llega hasta ahí, porque
las reglas concretas que gobiernan la forma de lo que la línea produce se
definen en la fase que las construye y las verifica.

Lo que este capítulo presenta es una síntesis del proceso y del argumento que
de él resultó, mientras que el detalle completo de cada instrumento, las
transcripciones de las entrevistas, la evaluación de cada solución frente a
cada eje y la ficha de cada requerimiento se conserva en los documentos de la
Fase 1, disponibles en la carpeta \texttt{Implementación/} del repositorio del
trabajo de
grado\footnote{\url{https://github.com/facundo1220/BackendServiceLine}}.

\section{Comprensión del problema y especificación}
\label{sec:comprension-problema}

La dimensión del problema de la fase consistió en aplicar un proceso de
obtención de requisitos sobre un caso de estudio, para comprender cómo se
inicia un proyecto backend y cómo se integran funcionalidades reutilizables en
la práctica real de un equipo de desarrollo.
El caso de estudio fue el Laboratorio Digital del Banco W
\citep{UniversidadIcesi2023}, cuyas prácticas, roles y restricciones
alimentaron la especificación de requerimientos de la herramienta.

\subsection{Caracterización del contexto de aplicación}
\label{sec:caracterizacion}

El Laboratorio Digital opera bajo marcos de trabajo ágiles, con Scrum y una
práctica de Dual Track que alinea el discovery con el delivery en sprints de
dos semanas \citep{MartyCagan2012}, y su composición de roles corresponde a la
de un equipo ágil, con desarrolladores, testers, un líder técnico y líderes de
producto. Su equipo técnico es reducido y estable, con cuatro desarrolladores
que atienden varios productos a la vez.

Como todo equipo que construye software dentro de una organización regulada,
trabaja en dos planos a la vez, porque conserva autonomía sobre la
construcción interna de sus soluciones, que es donde decide cómo se estructura
y se escribe el código de cada producto, y al mismo tiempo se subordina a las
definiciones de las funciones institucionales de tecnología, que establecen
los estándares a los que todo desarrollo debe ajustarse y las condiciones bajo
las cuales un servicio se integra con los sistemas centrales y llega a operar.
Esa frontera delimita sobre qué puede actuar legítimamente una herramienta de
generación, porque lo que produce pertenece por entero al primer plano y nada
de lo que genera puede alterar lo que la organización impone.

\subsection{Identificación y justificación de stakeholders}
\label{sec:stakeholders}

Definido el contexto, el paso siguiente fue establecer a quiénes consultar, y
la selección de roles condiciona lo que se va a encontrar, porque un
arquitecto, un desarrollador y un líder técnico ven cosas distintas del mismo
problema, y por eso el criterio se fijó antes de entrevistar. Se consideró
stakeholder a quien usa la herramienta, a quien fija reglas que esta debe
respetar, a quien puede habilitar o bloquear su adopción, y a quien ve cambiar
su trabajo por ella aunque no la use. Para ordenarlos se adoptó el modelo de
capas (\emph{onion model}) de Alexander \citep{Alexander2005}, que los agrupa
según su cercanía al artefacto, desde quienes lo operan hasta quienes solo
imponen condiciones sobre él.

La Tabla~\ref{tab:roles} recoge los roles consultados, la capa a la que
pertenecen y lo que se esperaba de la perspectiva de cada uno.

\begin{table}[H]
\centering
{\footnotesize
\begin{tabular}{p{3.2cm}p{2.6cm}p{9.4cm}}
\toprule
\textbf{Rol} & \textbf{Capa} & \textbf{Aporte esperado} \\
\midrule
Desarrolladores backend & Núcleo & Desagregación del proceso de arranque, puntos de dolor concretos y criterio sobre qué estandarizar \\
Testers & Núcleo & Requerimientos de consistencia externa de APIs e interfaces, y evidencia de heterogeneidad entre productos \\
Líder técnico & Sistema contenedor & Visión consolidada de la divergencia entre proyectos y consideraciones de gobernanza del catálogo \\
Líder de producto & Sistema contenedor & Impacto del tiempo de arranque e integración sobre la planeación de productos \\
Arquitectura & Entorno & Condiciones de diseño de un módulo de integración, organización de dominios y criterios de evolución \\
Desarrolladores externos & Entorno & Contraste sobre qué problemas son generales de la industria y qué mecanismos de reutilización aplican otras organizaciones \\
\bottomrule
\end{tabular}
}
\caption{Roles consultados, por capa del modelo de Alexander}
\label{tab:roles}
\end{table}

Los cuatro primeros roles pertenecen al equipo estudiado y los dos últimos no.
La función de arquitectura representa las definiciones a las que cualquier
desarrollo se subordina en una organización regulada, y fue la única fuente
que aportó condiciones de diseño en lugar de descripción de la práctica. Los
desarrolladores externos permiten separar lo propio del contexto estudiado de
lo que ocurre también en otras organizaciones. Las funciones consultadas son
arquetipos de cualquier equipo de producto con estas características, y es eso
lo que permite trasladar el análisis a otro contexto del mismo tipo.

\subsection{Protocolo de recolección de información}
\label{sec:protocolo}

La información se recolectó mediante entrevistas semiestructuradas de diseño
mixto, con nueve entrevistas repartidas en seis guiones, uno por cada rol
consultado. Una primera parte de preguntas abiertas recogió el proceso actual,
los puntos de dolor y el criterio de cada entrevistado, y es la que sostiene el
análisis cualitativo del que se derivan los temas, mientras que una segunda
parte más breve, de preguntas cerradas y respondida en los minutos finales,
produjo datos agregables que complementan las respuestas abiertas.

Esas preguntas cerradas se aplicaron antes de que existiera la primera línea de
código de la herramienta, de modo que son la única fuente de autoreporte
anterior al artefacto, y sobre ellas se constituyó la línea base con la que
compara la evaluación del Capítulo~\ref{cap:evaluacion}.

\subsection{Análisis cualitativo y temas derivados}
\label{sec:analisis-cualitativo}

Al ser las entrevistas de preguntas abiertas, cada entrevistado describió su
propia experiencia con sus propias palabras, de modo que el material se
transcribió y se analizó buscando qué se repetía de un testimonio a otro,
con el fin de encontrar los patrones de la problemática y no opiniones
aisladas. De ese análisis salieron quince temas, que recogen lo que los
entrevistados describieron de su práctica sin redactarlo como requerimiento ni
proponer una solución, y que se agruparon después en los cinco conjuntos de la
Tabla~\ref{tab:grupos}, porque es la agrupación la que permite analizarlos y la
que da origen a los requerimientos.

\begin{table}[H]
\centering
{\footnotesize
\begin{tabular}{p{2.9cm}p{1.7cm}p{6.9cm}p{3.1cm}}
\toprule
\textbf{Grupo} & \textbf{Temas} & \textbf{Lo que el testimonio muestra} & \textbf{Roles que lo evidencian} \\
\midrule
El costo del arranque & T-01, T-02, T-03 & Arrancar un proyecto es un proceso manual, concentrado en una persona, que impone un costo cognitivo inicial, se paga en alcance de negocio y resulta invisible para el negocio & Desarrolladores, líder técnico, líder de producto, externos \\
La reutilización y su propagación & T-04, T-05, T-09 & Existe un núcleo reconocido de funcionalidades que se repite entre productos, pero sin un mecanismo formal se recurre a copiarlo, y el contrato del servicio es el punto más frágil al propagar un cambio & Desarrolladores, arquitectura, líder de producto, externos \\
Causas estructurales de la divergencia & T-06, T-07, T-08 & Los proyectos divergen porque nada propaga las mejoras ni refuerza el estándar, porque se diseña para el alcance inmediato y porque el conocimiento transversal llega tarde & Líder de producto, líder técnico, arquitectura, externos \\
Condiciones de diseño & T-10, T-11, T-12, T-13 & La solución debe estandarizar con rigidez justificada, separar lo fijo de lo variable, evolucionar por versiones y poder verificarse & Todos los roles consultados \\
Condiciones de adopción y permanencia & T-14, T-15 & Adoptar requiere agilidad real y uso sencillo; permanecer requiere capacidad de evolución y un responsable definido & Desarrolladores, líder técnico, arquitectura, externos \\
\bottomrule
\end{tabular}
}
\caption{Los quince temas del análisis cualitativo, por grupo}
\label{tab:grupos}
\end{table}

El respaldo de los temas no es uniforme, porque catorce de los quince los
evidencia más de un rol, y tres de ellos aparecen en cuatro o cinco roles
distintos, lo que los sitúa como hallazgos compartidos entre perspectivas que
miran el mismo problema desde lugares diferentes. El único que sostiene un
solo rol es el que afirma que lo fijado en el código impide reutilizarlo
después, evidenciado por el líder de producto, y se conservó porque introduce
una causa estructural que ningún otro nombró, con el respaldo menor que eso
supone.

Los quince temas, agrupados así, son la unidad de trabajo del resto del
proyecto. Todo requerimiento posterior cita el tema o los temas que lo
originan, y el cierre del proyecto, en el Capítulo~\ref{cap:evaluacion}, se
formula contra los requerimientos y las brechas.

\subsection{Estudio comparativo de soluciones existentes}
\label{sec:estudio-comparativo}

Establecido qué muestra cada grupo temático, el paso siguiente fue medir qué
parte de eso resuelven ya las soluciones disponibles y qué parte no resuelve
ninguna, que es lo que determina si el problema planteado en la
sección~\ref{sec:planteamiento} queda cubierto por una herramienta existente.
La comparación se hizo sobre las ocho soluciones descritas en el estado del
arte (sección~\ref{sec:estado-del-arte}), que el Capítulo~\ref{cap:marco}
describe una por una y que esta sección evalúa frente a los pilares que los
temas establecieron, sin volver a describir ninguna.

Cada eje del marco de evaluación responde a un tema documentado con evidencia
de entrevistas, en lugar de a un inventario genérico de funcionalidades, de
modo que la comparación mide la distancia entre cada solución y el problema
real del contexto. Se excluyó el tema T-03, por tratarse de un efecto
organizacional y no de una capacidad evaluable en una herramienta, y los
catorce restantes quedan cubiertos por al menos un eje de la
Tabla~\ref{tab:ejes}, sin que ningún eje quede sin tema de origen.

\begin{table}[H]
\centering
{\footnotesize
\begin{tabular}{p{3.2cm}p{5.3cm}p{1.2cm}p{4.6cm}}
\toprule
\textbf{Enfoque} & \textbf{Eje} & \textbf{Temas} & \textbf{Qué significa cumplirlo} \\
\midrule
Generación de proyectos & Generación desde plantilla & T-01, T-02 & Producir un proyecto ejecutable a partir de una plantilla parametrizable, sin intervención manual \\
 & Contenido resuelto de origen & T-07, T-08 & Que el proyecto traiga ya resueltas las capacidades transversales y las restricciones que suelen descubrirse tarde \\
 & Externalización de la configuración & T-11 & Que lo susceptible de cambiar quede fuera del código y lo estable no \\
\midrule
Ciclo de vida & Composición de módulos en proyectos existentes & T-04, T-09 & Incorporar una capacidad transversal a un proyecto ya en curso, de forma desacoplada \\
 & Propagación de cambios de una capacidad compartida & T-05 & Que un cambio en una capacidad alcance a los proyectos que la usan sin intervenirlos uno por uno \\
 & Consistencia entre proyectos en el tiempo & T-06 & Detectar y corregir la desviación entre proyectos generados con versiones distintas \\
 & Versionado y distribución & T-12 & Que las capacidades y la herramienta evolucionen por versiones con adopción gradual \\
\midrule
Adopción y gobernanza & Adaptabilidad de lo generado & T-10 & Que el resultado se pueda modificar sin romper la relación con la herramienta \\
 & Verificabilidad & T-13 & Que lo generado y la herramienta puedan someterse a verificación automática \\
 & Costo de uso & T-14 & Que generar un proyecto cueste un comando o pocos comandos, sin configuración previa extensa \\
 & Gobernanza del catálogo & T-15 & Que exista un catálogo con propiedad, control de acceso y reglas de incorporación \\
\bottomrule
\end{tabular}
}
\caption{Los once ejes del marco de evaluación y el tema que cada uno atiende}
\label{tab:ejes}
\end{table}

La evaluación de las ocho soluciones frente a los once ejes, celda por celda,
se conserva en el documento \emph{Estudio comparativo de soluciones
existentes} de la Fase 1, disponible en el repositorio del trabajo de grado.
Su lectura agregada es que ninguna cubre los tres enfoques a la vez. Las
herramientas de plantilla resuelven la generación y el costo de uso, pero no
la composición sobre un proyecto en curso ni la gobernanza; las que sostienen
el ciclo de vida toman la plantilla completa como unidad de cambio; y las
plataformas de desarrollo interno resuelven el catálogo y la propiedad al
precio de desplegar y operar una aplicación propia.

\paragraph*{Brechas identificadas}
\label{sec:brechas}

Del cruce entre las ocho soluciones y los once ejes surgen tres capacidades
que el contexto necesita y que ninguna herramienta del mercado resuelve.

\textbf{B-01. Incorporar un módulo a un proyecto en curso.} Cuando meses
después un proyecto necesita autenticación, observabilidad o el consumo de un
servicio compartido, no hay forma de incorporarlo salvo a mano. La solución
más cercana modifica código existente en lugar de limitarse a crear archivos
nuevos, pero exige que todos los proyectos convivan en un mismo repositorio,
y el contexto estudiado trabaja con repositorios independientes.

\textbf{B-02. Propagar el cambio de un módulo.} Cuando una capacidad
compartida se copió a varios proyectos y después cambia, el cambio debe
llegar a todos sin intervenir cada uno a mano. Las soluciones más cercanas
registran la versión de plantilla que usó el proyecto y aplican los cambios
preservando las modificaciones locales, pero su unidad es la plantilla
completa. Lo que cambia en el caso documentado es una capacidad puntual y no
la estructura, de modo que actualizar la plantilla entera no propaga el
cambio de un módulo.

\textbf{B-03. Catálogo con reglas y bajo costo de uso.} Sostener un catálogo
exige propiedad definida, control de acceso y reglas de incorporación, sin
que eso encarezca el uso diario. Las soluciones más cercanas lo ofrecen, pero
exigen desplegar, configurar y operar una aplicación propia antes de generar
el primer proyecto, un costo desproporcionado para un equipo de cuatro
desarrolladores. En el extremo opuesto, las herramientas de un comando no
tienen catálogo alguno.

Las tres brechas cumplen una doble función, porque justifican construir algo
nuevo en lugar de adoptar una herramienta existente, y funcionan como criterio
de prioridad, porque un requerimiento que cierra una brecha es imprescindible
mientras que otro que solo mejora la experiencia de uso puede diferirse.

\subsection{Especificación de requerimientos}
\label{sec:especificacion-requerimientos}

De los quince temas, que expresan qué espera resolver el equipo, y de las tres
brechas, que delimitan lo que ninguna solución existente cubre, salieron
veintidós requerimientos, dieciséis funcionales y seis no funcionales. Con
ellos cierra el bloque de evidencia, porque son la forma verificable de todo
lo que las subsecciones anteriores establecieron.

A diferencia del tema, que describe lo que ocurre en la práctica, el
requerimiento fija un alcance concreto y establece qué debe hacer la
herramienta, sobre qué actúa y bajo qué condición se da por satisfecho. Por
eso cada uno se acompaña de su criterio de aceptación, que es el que vuelve
verificable aquello que el tema solo enunciaba, y la ficha de los veintidós,
con su enunciado, su prioridad, su origen y su criterio, se conserva en el
documento \emph{Especificación de requerimientos} de la Fase 1, disponible en
el repositorio del trabajo de grado.

\subsubsection*{Prioridad y cobertura}

Los requerimientos se priorizaron con la escala MoSCoW, cuyos niveles se
asignan según el origen de cada uno y no según la preferencia del
implementador. Es imprescindible el que responde a una de las tres brechas o
sin el cual la solución no resuelve el problema levantado; importante, el que
responde a un tema documentado pero cuya ausencia no invalida la
solución en esta fase; y deseable, el que aporta valor comprobado y puede
quedar fuera si el alcance lo exige.

Los dieciséis requerimientos funcionales se enuncian sobre las dos capacidades
centrales de la herramienta, generar proyectos e incorporar módulos, y sobre
las dos condiciones que las sostienen en el tiempo, la actualización y el
catálogo. La Tabla~\ref{tab:rf} recoge qué exige cada una y de qué temas y
brechas proviene.

\begin{table}[H]
\centering
{\footnotesize
\begin{tabular}{p{3.1cm}p{7.6cm}p{1.6cm}p{2.3cm}}
\toprule
\textbf{Capacidad} & \textbf{Qué exige} & \textbf{Req.} & \textbf{Origen} \\
\midrule
Generación de proyectos & Producir un proyecto ejecutable y parametrizable, con la estructura y el entorno de trabajo del marco ya resueltos, su documentación de API automática y un resultado determinista ante la misma entrada & RF-01 a RF-05 & T-01, T-02, T-06, T-07, T-08, T-11 \\
Composición de módulos & Incorporar un módulo a un proyecto existente sin que nadie edite archivos a mano, modificando lo que deba referenciarlo sin sobrescribir código propio, con consulta del estado y al menos un módulo real en el catálogo & RF-06 a RF-09 & T-04, T-09, T-15, B-01 \\
Actualización y propagación & Registrar en el proyecto qué módulos incorporó y en qué versión, actualizarlos sin intervención manual, informar de los desactualizados y preservar el código del equipo al hacerlo & RF-10 a RF-13 & T-05, T-06, T-10, B-02 \\
Catálogo y distribución & Declarar el catálogo de forma declarativa y versionada sin componentes de servidor, distribuir por el registro de paquetes de la organización y registrar la propiedad de cada activo & RF-14 a RF-16 & T-12, T-15, B-03 \\
\bottomrule
\end{tabular}
}
\caption{Los dieciséis requerimientos funcionales, por capacidad de la herramienta}
\label{tab:rf}
\end{table}

Los seis requerimientos no funcionales de la Tabla~\ref{tab:rnf} expresan los
atributos de calidad que la herramienta y lo que produce deben satisfacer,
con independencia de las operaciones concretas que ofrezca.

\begin{table}[H]
\centering
{\footnotesize
\begin{tabular}{p{1.4cm}p{2.9cm}p{8.9cm}p{1.3cm}}
\toprule
\textbf{ID} & \textbf{Atributo} & \textbf{Requerimiento} & \textbf{Origen} \\
\midrule
RNF-01 & Usabilidad & Las operaciones de la herramienta deben ejecutarse con un comando cada una & T-14 \\
RNF-02 & Adaptabilidad & El código generado debe ser modificable sin que ello impida seguir usando la herramienta sobre el proyecto & T-10 \\
RNF-03 & Verificabilidad & Lo generado debe poder someterse a análisis estático de código, incluidos aspectos de seguridad & T-13 \\
RNF-04 & Verificabilidad & La herramienta debe poder someterse a auditoría de dependencias y vulnerabilidades, con resultado consultable & T-13 \\
RNF-05 & Configurabilidad & El proyecto generado no debe contener valores fijos susceptibles de cambio, en particular direcciones de servicios y parámetros de negocio & T-11 \\
RNF-06 & Evolución & La herramienta y sus módulos deben evolucionar por versiones, admitiendo la convivencia de versiones distintas entre proyectos & T-12 \\
\bottomrule
\end{tabular}
}
\caption{Requerimientos no funcionales}
\label{tab:rnf}
\end{table}

Con la especificación queda cerrado el bloque de evidencia, y estos veintidós
requerimientos son el insumo del bloque siguiente, donde el marco de
referencia tendrá que habilitar cada uno o declarar a qué fase de
construcción se difiere. La comprobación de que cada uno quedó satisfecho se
realiza al cierre del proyecto, cuando existe el artefacto contra el cual
verificarla.

\section{Definición del marco de referencia}
\label{sec:marco-tecnico-solucion}

Establecida la evidencia del problema, el marco de referencia fija bajo qué
enfoque de ingeniería se construye la solución y sobre qué fundamentos se
apoya, de modo que las fases siguientes construyan dentro de límites ya
decididos en lugar de fijarlos durante la implementación. Lo que el marco
define no es la herramienta, sino la familia de proyectos que la herramienta
produce, porque es esa familia la que tiene que compartir una estructura para
que un mismo activo pueda servir a varios de sus productos.

Se desarrolla en dos partes, de las cuales la primera encuadra la
solución como una línea de producto de software, que es el enfoque que ordena
la construcción de los activos reutilizables y la de los productos que se
derivan de ellos, y declara con ese vocabulario qué productos forman la
familia y cómo se organiza su construcción. La segunda adopta los fundamentos
que dan forma al interior de cada producto y fija hasta dónde se toma cada
uno. El marco termina ahí, porque las reglas concretas que gobiernan la forma
de lo que la línea produce se definen y se verifican en la fase que las
construye, y se conservan en los documentos de esa fase.

\subsection{La solución como línea de producto de software}
\label{sec:caracterizacion-lps}

Una línea de producto de software es una familia de sistemas que comparten un
conjunto común y gestionado de características y que se construyen a partir de
activos comunes siguiendo un procedimiento establecido \citep{Clements2001}.
Lo que la separa de una colección de proyectos parecidos es que esas
similitudes están declaradas y gestionadas, de manera que un activo se
construye una vez y sirve a varios productos, en lugar de rehacerse en cada
uno. Encuadrar la solución bajo ese enfoque es la decisión de fondo del
trabajo, porque desplaza el problema desde elegir una herramienta hacia
definir las condiciones que gobiernan su uso, que es el plano en el que las
tres brechas pueden cerrarse para una familia entera de productos y no para un
proyecto a la vez.

\subsubsection*{Elementos constitutivos y pertinencia del enfoque}

El enfoque exige cuatro elementos constitutivos, y la ausencia de cualquiera
de ellos impide hablar de una línea de producto \citep{Clements2001}. Lo que
vuelve pertinente el encuadre no es que los cuatro puedan construirse, sino
que los cuatro ya estaban presentes en el contexto estudiado de forma
informal, cada uno con su evidencia en el análisis del problema. La
Tabla~\ref{tab:constitutivos} enfrenta lo que el enfoque exige con lo que el
contexto tenía y con lo que la línea construyó en su lugar.

\begin{table}[H]
\centering
{\footnotesize
\begin{tabular}{p{2.5cm}p{3.9cm}p{4.3cm}p{3.8cm}}
\toprule
\textbf{Elemento} & \textbf{Qué exige el enfoque} & \textbf{Cómo estaba en el contexto} & \textbf{Cómo se materializa en la línea} \\
\midrule
Alcance declarado & Declaración explícita de qué productos pertenecen a la familia & Una familia sin declarar, cuyos productos comparten propósito y estructura sin que nadie lo haya escrito & El alcance de la familia, con su condición mínima de pertenencia y sus exclusiones \\
Arquitectura común & Estructura compartida que hace posible que un mismo activo sirva a varios productos & Un patrón que el desarrollador debe dominar antes de ser productivo, transmitido imitando otros proyectos & Domain-Driven Design y arquitectura hexagonal, con una regla de dependencia verificable \\
Core assets & Activos reutilizables que incluyen arquitectura, contratos y modelos, y no solo código & Un núcleo reconocido de funcionalidades transversales que se repite entre productos y se copia & Las plantillas, los módulos, el \emph{contract package} y el catálogo \\
Mecanismo de derivación & Procedimiento que produce un producto a partir de los activos y de una selección de características & Copiar un proyecto y ajustarlo, que deriva sin trazabilidad ni vínculo con su origen & La herramienta de línea de comandos, construida en las fases posteriores \\
\bottomrule
\end{tabular}
}
\caption{Los cuatro elementos constitutivos de una línea de producto, frente a lo que el contexto ya tenía y a lo que la línea construyó}
\label{tab:constitutivos}
\end{table}

La tercera columna es la que decide, porque los cuatro elementos existían,
pero ninguno declarado ni gestionado, y esa es la diferencia que el enfoque
convierte en ventaja de ingeniería, porque lo que no está declarado no puede
compartirse de forma deliberada. La cuarta fila lo muestra en su forma más
clara, porque copiar un proyecto y ajustarlo es lo que la literatura denomina
\emph{clone-and-own}, una forma legítima de reutilización por su bajo costo
inicial, cuyo problema característico es la divergencia progresiva, porque
nada devuelve a las demás copias lo que se corrige en una y el costo crece con
el número de copias \citep{Dubinsky2013}. Es exactamente lo que el análisis
del problema encontró en campo, y lo que la primera brecha enuncia como
capacidad ausente.

De ahí se sigue el papel que le toca al artefacto, porque lo que el encuadre
aporta no es una herramienta más, sino las condiciones bajo las cuales se
emplean las que ya existen, y la que las fases siguientes construyen es el
mecanismo de derivación que esas condiciones exigen. Es además lo que vuelve
gestionados a los otros tres elementos, porque un alcance declarado y unos
activos comunes sin un procedimiento que los aplique vuelven a quedar en manos
de quien copia.

\subsubsection*{Organización general de la línea}

Los cuatro elementos describen qué debe existir, pero no cómo se articulan, y
para articularlos se adopta el modelo de referencia que cruza el espacio del
problema con el de la solución, y la ingeniería de dominio con la de
aplicación \citep{Apel2013}. Del cruce resultan cuatro cuadrantes, donde el
análisis del dominio declara el alcance de la familia y la variabilidad; la
implementación del dominio construye los activos comunes y la arquitectura de
referencia; el análisis de requerimientos de la aplicación es el momento en
que un desarrollador selecciona características para un producto concreto; y
la derivación del producto las ensambla. Los dos primeros corresponden a las
fases 1, 2 y 3 de este trabajo; los dos últimos ocurren cada vez que alguien
emplea la herramienta.

De los cinco flujos que conectan los cuadrantes, el que más condiciona el
diseño es el retorno de nuevos requerimientos hacia el dominio, porque es el
que permite que lo aprendido construyendo corrija el encuadre. El modelo
supone además un modelo de características formal sobre el cual se ejecuta la
selección, y este marco no lo define, de modo que la variabilidad de la línea
queda asignada a la Fase 2, que deberá formalizarla dentro de los límites que
el alcance de la familia y el catálogo de activos imponen.

\subsubsection*{Alcance de la familia de productos}

Pertenecen a la familia los servicios backend sobre Node.js y TypeScript que
responden por una sola capacidad de negocio, exponen al menos una API HTTP
documentada bajo el estándar OpenAPI y consumen servicios corporativos. Quedan
fuera las aplicaciones frontend, las que no exponen una API HTTP, los scripts
de propósito único y los servicios construidos sobre otros lenguajes o
ecosistemas.

Esa exposición es la condición mínima de pertenencia, y de ella depende que
todos los productos compartan la misma superficie de entrada, que es lo que
hace viable un activo común y lo que da sentido a la documentación OpenAPI
como capacidad transversal de la línea. Delimitar el alcance con esta
precisión lo exige el propio enfoque, porque uno demasiado amplio impide que
un activo sirva a todos los productos y uno demasiado estrecho no justifica el
esfuerzo de construir la línea frente al de construir un solo producto
\citep{Clements2002}.

La condición no fija un modelo de ejecución, porque un producto sin servidor
propio la cumple igual, y cuántas plantillas tiene la línea lo decide el
mecanismo de derivación y no el alcance de la familia. Que el marco absorbiera
ese hallazgo de la construcción sin alterar su condición de pertenencia es la
primera evidencia del flujo de retorno que el modelo de referencia declara.

\subsubsection*{Desviaciones declaradas respecto del modelo clásico}

El modelo adoptado opera sobre dos asunciones que el contexto no satisface, y
se declaran como desviaciones explícitas en lugar de forzar el caso dentro de
un esquema que no lo previó.

La primera es la de productos autónomos, porque el modelo no contempla retorno
desde el producto derivado, que una vez producido queda fuera del modelo,
mientras que aquí el proyecto se entrega a un equipo que conserva autonomía
para modificarlo. La línea no gobierna el producto después de derivado, sino
que mantiene con él un vínculo que permite propagar cambios sin imponer su
aceptación.

La segunda es la de derivación por etapas, porque en el modelo la selección de
características ocurre una sola vez y precede a la derivación, mientras que
aquí puede repetirse sobre un producto ya derivado, porque una capacidad
compartida debe poder incorporarse a un proyecto en curso.

Ambas se originan en la misma condición, que es que los productos pertenecen
a equipos autónomos con repositorios independientes. No son concesiones, sino
la consecuencia de aplicar el enfoque a un contexto que el modelo clásico no
previó.

\subsection{Fundamentos adoptados}
\label{sec:fundamentos-adoptados}

El encuadre anterior exige una arquitectura común y un modelo de dominio que
delimite y nombre lo que cada producto representa. Ambos se establecieron a
partir de los tres fundamentos de la Tabla~\ref{tab:fundamentos}, cuyas bases
conceptuales se presentaron en el marco teórico del Capítulo~\ref{cap:marco}.
Lo que aquí se fija es su adopción operativa, es decir qué se toma de cada uno,
para qué, y qué queda deliberadamente fuera.

\begin{table}[H]
\centering
{\footnotesize
\begin{tabular}{p{3.3cm}p{5.5cm}p{5.9cm}}
\toprule
\textbf{Fundamento} & \textbf{Qué aporta al marco} & \textbf{Hasta dónde se adopta} \\
\midrule
BIAN & El vocabulario de negocio y la unidad de delimitación de un producto & El Service Landscape como espacio de selección y el metadato de dominio; no gobierna el catálogo de capacidades técnicas \\
Domain-Driven Design & El modelado del interior del producto & Tres reglas del nivel estratégico y seis patrones tácticos; la relación entre \emph{bounded contexts} de proyectos distintos queda fuera \\
Arquitectura hexagonal & La separación entre el núcleo y todo lo que no es dominio & Tres principios, de los que se deriva una regla verificable por análisis estático \\
\bottomrule
\end{tabular}
}
\caption{Fundamentos adoptados y alcance de su adopción}
\label{tab:fundamentos}
\end{table}

\subsubsection*{Adopción de BIAN}

Se adopta la versión 14.0 del Service Landscape, y toda referencia a
\emph{Business Areas}, \emph{Business Domains} o \emph{Service Domains}
corresponde a esa versión. Declararla no es una formalidad, porque la
estructura del modelo cambia entre versiones y la 14.0 introdujo
renombramientos, eliminaciones e incorporaciones, particularmente en el área
de pagos.

BIAN gobierna tres aspectos del proyecto generado, que son la unidad de
derivación, el vocabulario con que se nombra lo que el producto representa y
el metadato de dominio que el producto registra. Igual de relevante es
delimitar lo que no gobierna, porque el catálogo de módulos agrupa
capacidades técnicas transversales y el estándar modela capacidades de
negocio bancario, no infraestructura técnica de aplicaciones. De los treinta
y ocho \emph{Business Domains} del Service Landscape, nueve corresponden al
tipo de productos que construye el contexto, y el subconjunto de
\emph{Service Domains} contenidos en ellos constituye el espacio de selección
por defecto del mecanismo de derivación.

\paragraph*{Unidad de derivación y criterio de granularidad}

Cada ejecución del mecanismo produce un proyecto que aloja exactamente un
\emph{bounded context}, correspondiente a un \emph{Service Domain}, y no se
admite derivar productos que alojen varios. La regla se apoya en el alcance
mínimo que define el estándar, porque un \emph{Service Domain} es una
capacidad de negocio discreta cuyo alcance es elemental, de modo que si se
descompone más, las partes pierden su contexto de negocio y pasan a ser
utilitarias. La propiedad de no solapamiento aporta la contraparte, porque
las capacidades que el estándar define no se solapan y la información que
gobiernan tampoco. Trasladada al contexto, impide que dos proyectos sean
dueños de la misma información.

\paragraph*{Registro del dominio como metadato}

Todo proyecto derivado registra, como metadato, el \emph{Service Domain} al
que corresponde y el \emph{Business Domain} que lo contiene. La regla admite
una excepción y establece una prohibición, de modo que si el producto
corresponde a una capacidad que no figura en el espacio de selección por
defecto, se permite declararla fuera de ese subconjunto registrando el
\emph{Business Domain} que la contiene. Lo que no se admite es omitir el
metadato, porque sin él la delimitación de dominio queda sin efecto
verificable y la alineación con el estándar se vuelve meramente declarativa.

\subsubsection*{Adopción de Domain-Driven Design}

Al corresponder cada proyecto a un único \emph{bounded context}, el nivel
estratégico se reduce a tres reglas, que fijan la correspondencia entre el
\emph{bounded context} y el \emph{Service Domain}, la obligación de un
lenguaje ubicuo derivado del estándar y la prohibición de alojar más de un
contexto por producto. La relación entre \emph{bounded contexts} de proyectos
distintos queda fuera del marco y se registra como punto de extensión.

En el nivel táctico se adoptan seis patrones, de los cuales cuatro son el
núcleo mínimo del modelado táctico, que son la entidad, el objeto de valor, el
agregado y el repositorio; los otros dos, el evento de dominio y el servicio
de aplicación, se incorporan por decisiones ya tomadas en el marco. Los
documentos del marco de referencia detallan cada uno y su realización en la
estructura del proyecto.

El término módulo exige una precisión, porque en este trabajo designa un
\emph{core asset} del catálogo. El patrón táctico homónimo no se adopta con ese
nombre, para evitar que un mismo término designe dos cosas distintas.

\subsubsection*{Adopción de la arquitectura hexagonal}

La arquitectura hexagonal aporta la regla que separa el núcleo del producto
de todo lo que no es dominio. De sus tres principios adoptados se deriva una
regla verificable que el marco toma como criterio, según la cual el núcleo no
importa ningún artefacto del exterior. Al construirse los productos en
TypeScript, esa regla puede comprobarse mediante análisis estático en lugar
de depender de revisión manual, lo que la convierte en una restricción
efectiva y no en una recomendación.

La articulación entre Domain-Driven Design y arquitectura hexagonal es
directa, porque los patrones tácticos describen qué contiene el núcleo, y la
arquitectura hexagonal describe cómo el núcleo se relaciona con lo que está
afuera. El repositorio es el punto donde ambos se encuentran, porque es un
patrón táctico cuya realización concreta vive en el exterior.

\subsubsection*{Aplicación conjunta de los tres fundamentos}

Los tres fundamentos operan de forma conjunta en cada derivación, y la
Tabla~\ref{tab:derivacion} recorre el camino completo, desde la selección en
el Service Landscape hasta el esqueleto generado, sobre una capacidad
concreta.

\begin{table}[H]
\centering
{\footnotesize
\begin{tabular}{p{3.6cm}p{4.6cm}p{6.5cm}}
\toprule
\textbf{Paso} & \textbf{Qué fundamento interviene} & \textbf{Qué produce} \\
\midrule
Selección de la capacidad & BIAN, Service Landscape & El \emph{Service Domain} y el \emph{Business Domain} que lo contiene \\
Delimitación del producto & BIAN y Domain-Driven Design & Un \emph{bounded context} único, con su lenguaje ubicuo \\
Modelado del núcleo & Domain-Driven Design, nivel táctico & La entidad raíz del agregado, sus objetos de valor, el contrato del repositorio y los eventos \\
Separación del exterior & Arquitectura hexagonal & Los puertos que el núcleo declara y los adaptadores que los realizan \\
Derivación del esqueleto & El mecanismo de la Fase 2 & La estructura de carpetas, los nombres y los contratos vacíos \\
\bottomrule
\end{tabular}
}
\caption{Recorrido de una derivación, desde el estándar hasta el proyecto generado}
\label{tab:derivacion}
\end{table}

El resultado sobre un producto de la familia es un núcleo en dos capas, la de
dominio dentro de la de aplicación, con los adaptadores que lo invocan desde
fuera por encima y los que él gobierna por debajo. La capacidad del estándar
no forma parte de esa arquitectura, sino que es lo que determina el contenido
del núcleo.

Lo que la derivación entrega es el punto de partida de un producto y no el
producto, porque las reglas de negocio las escribe después el equipo que lo
recibe, y esa distinción es la que mantiene al artefacto dentro de su
alcance.

\subsubsection*{Exclusión razonada de ISO 20022}

El estándar ISO 20022 se revisó como posible fundamento del marco y se
excluyó de forma deliberada. La razón es de naturaleza y no de dificultad, porque
ISO 20022 es un estándar de mensajería financiera que define modelos de
negocio y esquemas de mensajes para el intercambio de información entre
instituciones financieras, mientras que el artefacto de este trabajo no
produce, transforma ni transporta mensajes financieros, sino estructuras de
proyecto y componentes técnicos transversales. La relación entre ambos
estándares delimita el punto de contacto con precisión, porque BIAN
referencia correspondencias hacia ISO 20022 en la capa de especificación de
mensajes, y esa capa es exactamente la que el artefacto no produce.

La adopción de ISO 20022 corresponde, por tanto, a la lógica de negocio que
cada equipo implementa sobre el proyecto derivado, explícitamente fuera del
alcance de la solución. La exclusión se registra como punto de extensión y no
como omisión, porque un módulo de dominio que encapsule mensajería ISO 20022
sería un consumidor legítimo de este marco y no una parte de él.

Esta exclusión constituye una divergencia respecto del objetivo general
formulado en el anteproyecto, que comprometía la alineación con BIAN e ISO
20022. Se declara como tal, con su razón, y el Capítulo~\ref{cap:evaluacion}
la recoge entre las divergencias del proyecto al medir el cumplimiento de los
objetivos.

\section{Resumen del capítulo}

Este capítulo convirtió una intuición en un problema delimitado y un problema
en un encuadre. Al abrirlo existía la percepción de que arrancar un proyecto
cuesta demasiado y de que las funcionalidades transversales se copian; al
cerrarlo existen quince temas sostenidos por nueve entrevistas, tres
brechas que ninguna de las ocho soluciones estudiadas cubre, veintidós
requerimientos verificables y una arquitectura de la familia de productos que
la herramienta deberá materializar.

Tres decisiones estructurales sostienen el resultado, y la primera es el
encuadre como línea de producto de software, adoptado no por afinidad
conceptual sino porque los cuatro elementos que el enfoque exige ya estaban
presentes de forma informal, incluida la derivación por copia, cuya
limitación documentada es la que se encontró en campo. La segunda es la
delimitación de un producto por una capacidad de negocio del estándar, que da
a la línea una unidad de derivación y un vocabulario que no hay que inventar.
Y la tercera es que el marco diseña la familia de productos y no la
herramienta, lo que permitió verificar el encuadre antes de que existiera
código.

El capítulo dejó además dos límites declarados que el cierre del proyecto
recoge. La exclusión de ISO 20022 es una divergencia respecto del objetivo
general, y se declara con su razón en lugar de omitirse. Y las dos
desviaciones respecto del modelo de referencia clásico, el producto autónomo
y la derivación por etapas, no son concesiones sino la consecuencia de
aplicar el enfoque a un contexto de equipos autónomos con repositorios
independientes.

Lo que queda abierto por construcción es todo el mecanismo, porque el marco
acota lo que las fases siguientes pueden decidir pero no construye ni un
activo ni una línea de la herramienta. Con él quedan asignadas tres reglas
concretas que se definen donde se construyen: la forma que debe tener un
proyecto válido de la familia, los contratos por los que una capacidad
compartida se le acopla, y las condiciones que hacen válido a un módulo del
catálogo. Las dos primeras abren el Capítulo~\ref{cap:generacion}; la tercera
es del Capítulo~\ref{cap:gestor}.


%%----------------------------------------------------------
%%  CAPÍTULO 4: DISEÑO Y CONSTRUCCIÓN DEL MECANISMO DE GENERACIÓN DE PROYECTOS
%%----------------------------------------------------------

\chapter{Diseño y construcción del mecanismo de generación de proyectos}
%%==========================================================
%%  CAPÍTULO 4: DISEÑO Y CONSTRUCCIÓN DEL MECANISMO DE GENERACIÓN DE PROYECTOS
%%==========================================================
\label{cap:generacion}

El Capítulo~\ref{cap:fundamentacion} encuadró la solución como línea de
producto de software y fijó los fundamentos sobre los que se apoya. Un
encuadre, sin embargo, no reduce por sí solo el esfuerzo de iniciar un
proyecto; lo reduce un mecanismo que lo aplique sin intervención manual. Esta
fase construyó ese mecanismo, y con él las reglas concretas que el marco dejó
asignadas, que son qué forma tiene un proyecto válido de la familia y por qué
contratos se le acopla una capacidad compartida.

Lo que entregó es una herramienta de línea de comandos instalable que, a
partir de los activos de la línea y de una selección de características,
produce un proyecto backend conforme al marco. La condición de llegada se
formuló antes de construir y es verificable, porque exige un producto derivado
que arranca sin que nadie edite un archivo, que supera su propio análisis
estático con la regla de dependencia activa como error, y que declara con qué
versión de cada activo fue producido. El detalle de la construcción, con sus
decisiones intermedias y la evidencia de cada verificación, se conserva en los
documentos de la Fase 2, disponibles en el repositorio del proyecto.

\section{Hitos de la construcción}
\label{sec:hitos-fase-2}

La fase se ejecutó en cinco hitos, y el orden entre ellos es en sí mismo una
decisión, porque los que construyen los activos van antes que el que construye
el generador. Todo lo que un activo contiene puede verificarse a mano mientras
no exista todavía ningún generador, de modo que invertir el orden habría
obligado a descubrir un error de la estructura cuando ya existía una
herramienta edificada alrededor de ella. La Tabla~\ref{tab:hitos-generacion}
recoge qué dejó cada hito y cómo quedó comprobado.

\begin{table}[H]
\centering
{\footnotesize
\begin{tabular}{p{2.9cm}p{6.1cm}p{5.4cm}}
\toprule
\textbf{Hito} & \textbf{Qué dejó construido} & \textbf{Cómo quedó comprobado} \\
\midrule
La forma del producto & Las capas, la regla de dependencia, la frontera del contrato externo y los contratos de integración de la línea & El análisis estático falla nombrando la regla cuando se introduce una importación prohibida \\
La plantilla base & Un proyecto ejecutable con el entorno de trabajo configurado y la documentación de API automática & Arranca sin que nadie edite un archivo y supera su propia verificación \\
El catálogo & El manifiesto declarativo, el modelo de gobierno y la distribución por el registro de la organización & Se publica desde su repositorio y se consume en ejecución, sin plataforma propia \\
La herramienta & El núcleo común con sus seis puertos, el punto de extensión y la operación de derivación & La suite corre sin tocar disco, red ni terminal, con dobles de cada puerto \\
El almacén de plantillas & Una plantilla que se describe a sí misma, el motor de gramática cerrada y la forma del producto como contrato & Una segunda plantilla, de otro modelo de ejecución, entró sin que cambiara una línea de la herramienta \\
\bottomrule
\end{tabular}
}
\caption{Hitos de la construcción del mecanismo de generación}
\label{tab:hitos-generacion}
\end{table}

El quinto hito no estaba previsto al abrir la fase, y se abrió al comprobar que
el alcance comprometía un almacén de plantillas, en plural, y que la línea
tenía una sola plantilla sin lugar donde poner una segunda; la
sección~\ref{sec:almacen} narra qué obligó a revisar y qué quedó más firme
que antes.

\subsection{Alcance de la fase}
\label{sec:alcance-fase-2}

La frontera con la fase siguiente no es temporal sino funcional, y esa
distinción evita el riesgo principal del diseño, porque ambas fases construyen
la misma herramienta. Lo que toda operación comparte se construye aquí, y lo que
es propio de una operación se construye en la fase que la introduce. En
particular, el punto de extensión por el que se agregan operaciones nuevas
tiene que existir antes de que la Fase 3 agregue las suyas; de lo contrario
esa fase rehace la herramienta en lugar de ampliarla. La
Tabla~\ref{tab:alcance-fase-2} recoge lo que queda fuera y por qué.

\begin{table}[H]
\centering
{\footnotesize
\begin{tabular}{p{6.2cm}p{8.8cm}}
\toprule
\textbf{Queda fuera} & \textbf{Razón} \\
\midrule
Incorporación de módulos a un producto & El marco trata generar y componer como operaciones distintas sobre el mismo mecanismo \\
Construcción de módulos del catálogo & La composición se valida en la fase donde el módulo tiene función \\
Actualización sobre productos existentes & Requiere más de una versión publicada de un activo consumido \\
Despliegue e infraestructura & Excluidos en la especificación de requerimientos \\
\bottomrule
\end{tabular}
}
\caption{Alcance diferido de la fase}
\label{tab:alcance-fase-2}
\end{table}

Todo lo diferido se registra como punto de extensión y no como limitación de
la solución, que es la diferencia entre una frontera decidida y una carencia.

\section{La forma de un producto de la familia}
\label{sec:organizacion-codigo}

Los fundamentos adoptados establecen qué contiene el núcleo de un producto y
cómo se relaciona con el exterior, pero no dónde reside físicamente cada
elemento. Eso lo fijó la línea antes que ninguna otra cosa, porque sin una
forma declarada no hay nada que derivar.

El proyecto se organiza primero por capa y, dentro de la capa de dominio, por
agregado.

\subsection{Las reglas de dependencia, y por qué son verificables}

Las dependencias entre capas admiten una sola dirección, porque la
infraestructura puede depender de la aplicación y del dominio; la aplicación,
únicamente del dominio; y el dominio, de ninguna de las dos. La restricción no
alcanza a dos elementos, que son el punto de composición, donde se asocian los
puertos con sus adaptadores, y el contrato de la línea, que solo declara
tipos.

De ahí la regla que la línea adopta como criterio: ningún elemento de dominio
importa un elemento de aplicación o de infraestructura. El dominio no conoce
el framework, el motor de persistencia ni el protocolo de transporte.

Lo que convierte esa regla en una restricción efectiva no es enunciarla sino
cómo se comprueba. Se verifica por análisis estático sobre las rutas de
importación, no por revisión manual, y la violación se reporta como error y
no como advertencia, de modo que un proyecto que la incumpla no supera su
verificación. Una regla que solo advierte no restringe.

Las convenciones de nomenclatura que acompañan a esta estructura, con el
sufijo que corresponde a cada patrón táctico y la forma de derivar un nombre
desde el estándar, se recogen en los documentos de la Fase 2.

\subsection{La frontera del contrato externo y la documentación de API}

Los tipos que un adaptador de entrada expone al exterior son distintos de los
objetos de dominio, y el adaptador traduce explícitamente entre unos y otros,
porque son dos contratos con ciclos de vida independientes. El objeto de
dominio cambia cuando cambian las reglas del negocio; el tipo expuesto, solo
cuando cambia el acuerdo con quienes consumen el servicio. Si el adaptador
expusiera el objeto de dominio, cualquier reorganización interna se
convertiría en un cambio incompatible para los consumidores.

Esa separación tiene una consecuencia aprovechable que cierra un
requerimiento, porque los tipos expuestos, al ser explícitos, son la fuente
desde la cual se deriva la documentación de API del producto, porque el
documento se arma con los esquemas que el producto ya usa para validar sus
peticiones y respuestas, y no con una declaración paralela. Es lo que impide
que la documentación y el comportamiento se separen, porque no hay dos fuentes
que mantener sincronizadas sino una. El documento lleva además, para cada
operación, el identificador y el grupo que el estándar le da, que es lo único
que permite leerlo contra el publicado cuando la forma de las rutas del
producto puede variar.

La interfaz de exploración se decidió dos veces, y la segunda decisión no
contradijo las razones de la primera. Al construir el documento se descartó un
visor del tipo \emph{Swagger UI}, porque exigiría una dependencia de ejecución
en todo producto o un script traído de internet, y ninguna de las dos es
admisible en el entorno al que la línea sirve. Lo que esa decisión dejaba a
cada equipo, la línea lo resolvió después una sola vez, de modo que el
producto sirve una ruta que renderiza el mismo documento con el visor cargado
de la copia instalada con el producto y nunca de internet, con la dependencia
declarada como de desarrollo. Así no cambia lo que un servicio desplegado
carga, la auditoría de dependencias de ejecución sigue dando el mismo
resultado, y una variable de entorno decide si la ruta existe, apagada por
defecto en producción.

\section{Puertos y contratos de integración}
\label{sec:puertos-contratos}

La estructura admite que la capa de aplicación dependa de contratos de
puerto, pero no dice qué son ni quién es su dueño. Resolverlo es la condición
de todo el mecanismo de composición, porque de ahí depende cómo un activo
reutilizable se acopla a un producto sin romper el desacoplamiento que la
regla de dependencia establece. Es también el punto donde la arquitectura
hexagonal se sostiene o se vuelve nominal, porque si el módulo declarara el
puerto, el producto acabaría dependiendo del módulo y la dirección quedaría
invertida en la práctica, aunque la definición afirmara lo contrario.

\subsection{El contract package}

Un módulo debe poder servir a más de un producto, y para ello alguien tiene
que declarar el puerto. La solución es un artefacto neutral del que dependen
tanto los productos como los módulos, y del que ninguno de los dos depende
del otro. Es el \emph{contract package}, un \emph{core asset} que contiene
las declaraciones de los puertos técnicos de la línea.

Lo que hace viable esta decisión es que la dependencia de un producto hacia él
es exclusivamente de tipos, y una importación de tipos se elimina durante la
compilación, de modo que el producto declara el contrato pero en ejecución no
existe referencia alguna a él. Sin esa propiedad sería una dependencia real y
la solución perdería la mitad de su valor.

De ahí su regla de contenido, que es la ausencia de todo aquello que
sobreviva a la compilación. No se enuncia como «solo interfaces», dado que
una interfaz necesita tipos para sus parámetros y retornos. Cuando un puerto
opera sobre datos del dominio de un producto se declara genérico y el
producto aporta el tipo, con lo que el contrato nunca conoce el dominio de
ningún producto y el producto nunca depende del módulo.

\subsection{Las dos clases de puerto}

El criterio que separa las dos clases de la Tabla~\ref{tab:puertos} es la
firma y no el uso: un puerto es de dominio si su firma menciona conceptos del
dominio del producto, y es técnico si su firma es expresable sin conocer ese
dominio.

\begin{table}[H]
\centering
{\footnotesize
\begin{tabular}{p{2.4cm}p{4.8cm}p{3.5cm}p{4.0cm}}
\toprule
\textbf{Clase} & \textbf{Criterio} & \textbf{Quién lo declara} & \textbf{Quién lo realiza} \\
\midrule
De dominio & Su firma menciona conceptos del dominio del producto, como el repositorio de un agregado & El producto, en su capa de dominio & El propio producto, en su infraestructura \\
Técnico & Su firma es expresable sin conocer el dominio, como el registro de eventos o el cifrado & El \emph{contract package} & Un módulo del catálogo \\
\bottomrule
\end{tabular}
}
\caption{Las dos clases de puerto y la propiedad de cada una}
\label{tab:puertos}
\end{table}

La consecuencia es que ningún módulo puede apropiarse de un puerto de
dominio. Un repositorio expresa intención de negocio en el lenguaje del
producto, y no puede provenir de un activo genérico sin arrastrar el dominio
hacia el catálogo. En sentido inverso, un módulo sí puede proponer un puerto
técnico nuevo, porque exigir aprobación previa para cada uno convertiría el
contrato en el cuello de botella del catálogo. El caso que obliga a ejercitar
las dos clases a la vez es la persistencia, y se desarrolla en el
Capítulo~\ref{cap:gestor}, donde el módulo que la realiza se construye.

\section{La variabilidad de la línea}
\label{sec:variabilidad}

El modelo de variabilidad declara qué puede diferir entre un producto y otro.
Es el contrato que la herramienta consume, y por eso se escribe en un formato
declarativo legible por máquina, porque sin él la parametrización se decidiría
sobre la marcha durante la implementación y dejaría de ser verificable.
Reside en artefactos versionados, cada uno con su esquema, que la
verificación de la línea valida en cada entrega.

Las decisiones se agrupan en las tres clases de la
Tabla~\ref{tab:clases-variabilidad}, y la clase determina la forma de la
respuesta.

\begin{table}[H]
\centering
\begin{tabular}{p{2.6cm}p{5.4cm}p{6.6cm}}
\toprule
\textbf{Clase} & \textbf{Qué resuelve} & \textbf{Forma de la decisión} \\
\midrule
Identidad & Qué representa el producto y cómo se nombra & Un valor escrito o elegido \\
Selección & Qué capacidades del dominio incorpora & Un subconjunto de un conjunto declarado \\
Presencia & Qué piezas opcionales se incluyen & Inclusión u omisión \\
\bottomrule
\end{tabular}
\caption{Clases de variabilidad del modelo}
\label{tab:clases-variabilidad}
\end{table}

La clase de selección es la que sitúa este modelo dentro del cuadrante de
análisis de requerimientos de la aplicación, porque es el punto donde el
desarrollador elige características y no solo nombra lo que se genera. El
modelo declara trece puntos de variación, cada uno con su clase, su valor por
defecto y las condiciones que lo habilitan o lo excluyen.

\subsection{El espacio de selección de dominio}

El \emph{Service Domain} no se escribe libremente, sino que se elige de un
conjunto declarado. La escritura libre dejaría el metadato de dominio sin
validación posible y convertiría la alineación con el estándar en una
afirmación no verificable, que es precisamente lo que el marco buscaba
evitar. Cada entrada describe una capacidad de negocio con su área, su grupo,
la entidad de información que gobierna y las operaciones que el estándar
publica sobre ella. El estándar modela cada capacidad en dos niveles y el
espacio recoge ambos, de modo que al derivar se ofrecen tanto las operaciones
que gobiernan el ciclo de vida de la capacidad como las de sus subdivisiones.

Una entrada existe cuando sirve a un propósito escrito del proyecto, y se
agrega junto con ese propósito. No se agrega para poblar un catálogo ni para
cubrir un \emph{Business Domain}, porque transcribir entero el subconjunto
que el marco declaró convertiría una lista de referencia en una asignación
definitiva y pondría a elegir entre capacidades que el contexto no va a
construir.

\paragraph*{La fuente del dato, y su límite}

El dato de una capacidad se transcribe del documento que el estándar publica
para ese \emph{Service Domain}, y no de una vista general del
\emph{Service Landscape} ni de una paráfrasis. La distinción no es formal, y lo
demostró una transcripción hecha sin ese documento, que atribuyó a una
capacidad una cláusula que pertenece a otra, le asignó una operación en un
nivel que el estándar no usa y declaró como no mutante una operación que el
documento
publica con cuerpo de petición. Los tres errores se corrigieron al conseguir
la fuente.

La consecuencia es un límite declarado y no un hueco oculto, porque solo las
capacidades cuyo documento publicado está a la vista se transcriben, y las
demás conservan su descripción parafraseada hasta que lo esté. Inventar el
contenido del estándar sería el modo de perder exactamente la propiedad que
justifica adoptarlo.

\paragraph*{La capacidad que el catálogo no ofrece}

El mecanismo admite declarar un \emph{Service Domain} fuera del espacio por
defecto, registrando el \emph{Business Domain} que lo contiene, y prohíbe
únicamente omitir el metadato. El contexto lo necesita porque no todo backend
del laboratorio corresponde a una capacidad que el estándar publique.

Lo que se mantiene cerrado es la \emph{forma} de la respuesta, no siempre el
conjunto de valores. Una capacidad declarada elige su \emph{Business Domain}
entre los que el marco reconoce y sus operaciones entre los términos de
acción que el mecanismo deriva, y el \emph{Business Area} se lee del grupo en
lugar de preguntarse, lo que impide que un producto responda las dos cosas
por separado y se contradiga con el estándar. El metadato sigue validado, de
modo que lo que cae es la forma y no la razón.

El registro del producto marca el caso, y esa marca habilita una pregunta que
antes no tenía dónde formularse: cuántos productos de la línea responden por
una capacidad que el estándar no publica, y cuáles de ellas se repiten. Tres
productos declarando la misma capacidad es la evidencia que justifica
incorporarla al catálogo, y constituye el flujo de retorno de requerimientos
hacia el dominio que el modelo de referencia declara.

\section{De la plantilla base al almacén de plantillas}
\label{sec:almacen}

\subsection{La plantilla base como validación del marco}
\label{sec:activo-base}

La plantilla base es el proyecto del que se derivan todos los demás, y el
primer artefacto ejecutable de este trabajo. Materializa la estructura de la
sección~\ref{sec:organizacion-codigo} y trae el entorno de trabajo ya
configurado, con el análisis estático y su regla de dependencia activada como
error, el ejecutor de pruebas, las herramientas de formato y la configuración
externalizada que RNF-05 exige.

Su función va más allá de servir de punto de partida, porque es la validación
del marco. Si no arranca, o no supera su propio análisis estático, el problema
está en la definición y no en el código, y hay que corregir la definición
antes de seguir. Esa es la razón por la que se construye antes que la
herramienta.

Una decisión de este hito condiciona el resto de la fase, y es que la plantilla
base no se mantiene a mano, sino que es exactamente el resultado de rendir las
plantillas con una entrada determinada. La verificación lo comprueba archivo
por archivo, con la única excepción de los tres que llevan la identidad del
artefacto, que en el activo describen un activo y en un producto describen un
producto. La consecuencia es que no pueden divergir, porque un cambio en una
plantilla que no se refleje en el activo rompe la verificación y un retoque
manual del activo la rompe igual. Es el mecanismo que impide que la plantilla
y su referencia se separen en silencio, que es el modo habitual en que un
generador deja de producir lo que su ejemplo muestra.

\subsection{Una plantilla es un activo que se describe a sí mismo}

Con una sola plantilla, el reparto entre lo que la plantilla declara y lo que
la herramienta sabe no tiene que decidirse, porque todo puede vivir en la
herramienta sin que se note. El alcance del proyecto compromete otra cosa, ya
que habla de un sistema de almacenamiento de plantillas y RF-14 del catálogo
de plantillas y módulos, en plural. Hacer exigible esa promesa obligó a hacer
el reparto explícito, y lo que salió es más firme que lo que había.

Una plantilla pasa a ser un activo del catálogo del mismo rango que un
módulo, y como tal se describe a sí misma en un manifiesto que viaja en su
paquete, declarando el contrato contra el que se construyó, qué forma de
producto deriva, dónde viven sus archivos y de qué está hecho su producto. El
manifiesto va en el paquete y no en el catálogo por el mismo argumento que
rige para un módulo, porque lo que una plantilla declara de sí misma cambia
con ella, y ponerlo en el catálogo obligaría a coordinar dos repositorios para
publicar una versión. Sus archivos salieron de la herramienta al paquete de
la plantilla, con lo que una plantilla de otro equipo puede vivir en el
repositorio de ese equipo y una versión de la plantilla deja de costar una
versión de la herramienta.

\paragraph*{La forma del producto, como contrato}

La línea había descrito en prosa los puntos donde un módulo inserta su
cableado, y la herramienta los conocía como constantes de su propio código.
Con una plantilla eso bastaba; con dos deja de bastar, porque una plantilla
que ponga su punto de composición en otro archivo sería inservible. La forma
del producto pasó entonces a ser un dato con nombre y versión, declarado en
tres lugares que se comparan entre sí: la plantilla declara qué forma derivan
sus productos, el producto la lleva escrita en su registro, y el módulo
declara a qué formas sirve.

Un punto de inserción de los ocho que la forma hexagonal declara no dice solo
dónde se inserta, sino con qué sintaxis se comenta el archivo que lo contiene,
que es lo que permite que el mismo motor escriba en un archivo de TypeScript y
en uno de configuración sin conocer ninguno de los dos.

El producto lleva la forma y no se lee de la plantilla, por la misma razón
por la que registra con qué versión del activo fue derivado. La
compatibilidad de un módulo se decide contra el producto, y la plantilla que
lo derivó puede no estar instalada, no estar publicada o no existir. Lo que
el producto es tiene que poder leerse del producto.

El cambio tuvo además una consecuencia de la clase que una definición en prosa
no obliga a resolver. El marco suponía que
el producto había derivado del activo que escribió esas marcas, y nadie lo
comprobaba. Hoy son dos rechazos de la composición: uno cuando un módulo
reemplaza una región que la forma del producto no declara, y otro cuando al
producto le falta un punto que su propia forma declara. Un supuesto implícito
pasó a ser comprobable por los dos lados.

\subsection{El motor y la frontera de lo expresable}

Con un solo activo base el motor podía limitarse a sustituir nombres, y la
propiedad que lo protegía era que las plantillas no deciden nada. Un almacén
con más de una plantilla obliga a revisar esa propiedad, porque una plantilla
que no pueda expresar una condición tiene que duplicarse en dos archivos por
cada decisión, y eso fue lo que llevó a que una sola plantilla ocupara más de
cincuenta archivos.

El motor ganó entonces condición y repetición, pero con una gramática
cerrada, que es una propiedad más fuerte que la anterior y no una renuncia.
Antes se afirmaba que las plantillas no deciden nada; ahora se afirma que
todo lo que una plantilla puede decidir está enumerado, en tres
construcciones y dos tablas cerradas. No hay evaluación dinámica de código,
no existe forma de que una plantilla nombre una función, y la condición la
lee el mismo evaluador que ya leía las restricciones del modelo de
variabilidad, de modo que no hay una segunda gramática que mantener. Una
etiqueta que el motor no reconoce es un defecto de la plantilla y aparece en
la verificación del activo que la posee. Con esa capacidad, los archivos de
plantilla se redujeron casi a la mitad.

La frontera de lo expresable se escribió como parte de la decisión, y dos
cosas se eliminaron en lugar de expresarse. El formato por ancho de línea
salió porque es formato y lo comprueba el verificador del propio producto, y
el impresor de literales del documento de API pasó a ser texto en su
plantilla. Lo que queda del lado de la herramienta queda por una razón
declarada, porque es mecanismo y no producto: las huellas de integridad, el
reloj, el respaldo, las preguntas, el registro, el catálogo y los códigos de
salida.

El efecto sobre el código fue que la función que escribía el producto dejó de
ser el catálogo de archivos del producto escrito en código y pasó a ser un
intérprete, con una firma que no nombra ni el framework, ni el estándar, ni
ninguna carpeta de un producto.

\subsection{La evidencia de que el mecanismo es genérico}

Todo lo anterior descansa en una proposición: que lo que queda en la
herramienta es un mecanismo y no la descripción de un ejemplo. Con una sola
plantilla la proposición no es comprobable, porque un intérprete genérico y
una función escrita para ese caso se comportan igual.

La evidencia se produjo en dos formas que prueban cosas distintas, y la primera
es una plantilla de otra familia como fixture de prueba, en otro lenguaje,
con otro estilo de comentario, un trabajo por lotes en lugar de un servicio y
un contrato que no es el de esta línea; prueba que el intérprete deriva una
familia de entrada distinta. La segunda es una plantilla publicable que
deriva un API HTTP sin servidor propio, con otra forma de producto; prueba lo
que el fixture no puede, que es que una plantilla de otra forma pasa las
puertas de la línea, deriva un producto que registra esa plantilla y esa
forma, se instala y supera su propia verificación.

Lo que la segunda plantilla difiere no es un segundo caso de lo mismo. No
tiene contenedor de inyección, no tiene servidor propio, y sus puntos de
inserción son cinco en lugar de ocho, con uno de los roles que sencillamente
no existe. Ese último dato es el hallazgo del contrato de forma, porque el
archivo de configuración de esa plataforma es JSON estricto y un punto de
inserción es una línea de comentario que ahí dejaría el archivo inválido. Una
forma cuyo archivo de configuración es JSON no puede tener punto de entorno, y
el contrato lo expresa sin inventar nada, porque la forma declara los puntos
que tiene y el mecanismo rechaza nombrando el que falta.

La ronda dejó además probado lo que era una condición de diseño declarada y
nunca ejercida, que es que el diseño no quedó atado a un framework concreto.
Una plantilla sobre otro modelo de ejecución entró sin que una línea de la
herramienta cambiara.

\section{Gobierno de los activos y catálogo}
\label{sec:gobierno-activos}

Una línea de producto sostiene tres actividades en paralelo: el desarrollo de
los activos reutilizables, el de los productos que se derivan de ellos, y la
gestión que decide el alcance, asigna responsables y sostiene la línea en el
tiempo \citep{Clements2001}. Las dos primeras quedaron resueltas con los
activos y el mecanismo; la tercera es la que responde a la brecha B-03, porque
sin gestión los activos derivan, nadie responde por ellos y el catálogo se
convierte en una lista de cosas que alguien publicó alguna vez.

\subsection{Tipos de activo y escala de adaptación}

El catálogo reconoce cuatro tipos de \emph{core asset}, de los que la
plantilla aporta la estructura de partida de un producto; el módulo, la
realización de una o más capacidades técnicas; el \emph{contract package}, las
declaraciones de los puertos técnicos de la línea; y el manifiesto del
catálogo describe a los demás. Tres condiciones son comunes a todos: un
responsable declarado, un versionado que distinga el cambio incompatible del
compatible, y la posibilidad de saber qué productos lo consumen. La tercera es
la que hace gobernable al conjunto, porque permite responder qué se ve
afectado por un cambio consultando el catálogo, sin abrir los repositorios de
los productos. El responsable es un equipo y no una persona, porque un activo
cuyo responsable fuera una persona quedaría huérfano al salir esa persona de
la organización. Las condiciones específicas por tipo las detalla el documento
de gobierno de los activos de la Fase 2.

El equipo que recibe el producto derivado conserva la autonomía de
modificarlo y, al mismo tiempo, la línea debe poder propagarle cambios. La
escala de adaptación de la Tabla~\ref{tab:adaptacion} resuelve esa tensión
declarando tres niveles con consecuencias distintas.

\begin{table}[H]
\centering
{\footnotesize
\begin{tabular}{p{2.8cm}p{5.6cm}p{6.3cm}}
\toprule
\textbf{Nivel} & \textbf{Qué hace el equipo} & \textbf{Efecto sobre la propagación} \\
\midrule
Configuración & Ajusta valores previstos por el activo, fuera del código & El producto sigue recibiendo toda mejora del activo \\
Reemplazo & Sustituye el adaptador de un puerto por otro propio, conservando el contrato & El producto sigue recibiendo las mejoras que no tocan el adaptador reemplazado \\
Copia desvinculada & Copia el código del activo al producto y lo modifica & El producto deja de recibir mejoras, y el catálogo lo registra como desvinculado \\
\bottomrule
\end{tabular}
}
\caption{Escala de adaptación y su efecto sobre la propagación de cambios}
\label{tab:adaptacion}
\end{table}

Los dos primeros niveles convierten la adaptabilidad en una propiedad
arquitectónica y no en una concesión. El tercero se admite en lugar de
prohibirse por una razón práctica, porque prohibirlo no evitaría que ocurra,
solo lo dejaría sin registro. Una copia desvinculada se lee entonces como
señal de que al activo le falta un punto de configuración o de reemplazo, y
alimenta el flujo de nuevos requerimientos hacia el análisis del dominio.

\subsection{Versionado y cambio incompatible}

La línea fija qué constituye un cambio mayor, uno menor y un parche, porque
de ello depende que dos productos puedan operar con versiones distintas de un
mismo activo, que es el criterio de RNF-06. El documento de gobierno de los
activos de la Fase 2 enumera cada caso.

Lo que importa es el criterio con que se decide, porque la compatibilidad se
determina por inspección de lo declarado, es decir la firma de las
operaciones, el esquema de configuración y el identificador de registro. Toda
la evidencia necesaria está en el contrato, sin recurrir a la ejecución, y
con el lenguaje adoptado el compilador detecta por sí mismo los cambios de
firma que rompen a un consumidor. De ello se sigue una limitación asumida de
forma explícita, y es que el activo no se comprueba ejecutándolo, de modo que
un cambio de comportamiento que no altere la firma pasa inadvertido. Se registra
como punto de extensión.

El \emph{contract package} se organiza en submódulos por capacidad pero se
versiona como una unidad, de modo que un cambio incompatible en una capacidad
eleva la versión mayor de todas. La consecuencia se asume, porque un producto
que consume una sola capacidad recibe una versión mayor que no le afecta, y
un módulo puede verse obligado a republicarse sin que su implementación haya
cambiado. La separación en artefactos independientes queda prevista, con una
señal de cuándo conviene hacerla: cuando un módulo deba republicarse sin
alterar su implementación más de dos veces, el acoplamiento de versión cuesta
más de lo que ahorra en gobernanza.

\subsection{El catálogo, su gobernanza y la resolución en ejecución}
\label{sec:catalogo-operacion}

El catálogo es un manifiesto declarativo y versionado que reside en un
repositorio propio y que \emph{describe} activos sin contenerlos: una entrada
nombra el paquete y su responsable, de modo que un módulo construido por otro
equipo vive en el repositorio de ese equipo y el catálogo solo lo registra.
Esa forma es lo que satisface RF-14 sin desplegar ni operar ninguna
plataforma, y lo que cierra la brecha B-03 al alcance de un equipo pequeño.
Los campos de una entrada y los estados por los que pasa se recogen en ese
mismo documento.

Incorporar una entrada exige que el activo satisfaga las condiciones de
validez y que la propuesta sea revisada por un rol responsable del gobierno
de la línea, cuya existencia define el marco mientras que a quién se le
asigna lo decide cada organización que lo adopte. No se admite retirar un
activo, porque uno del que dependen productos existentes no puede eliminarse
sin romperlos. El estado de desuso es terminal, el activo deja de ofrecerse al
derivar y permanece en el catálogo indefinidamente. El retiro acompañado
de un procedimiento de migración se registra como trabajo futuro.

La forma en que la herramienta obtiene el catálogo cambió al medir su costo, y
el episodio ilustra un criterio. La primera versión fijaba la versión del
catálogo dentro de la herramienta, con el argumento de que RF-05 exige que la
misma entrada produzca el mismo resultado. Medido, ese argumento no se
sostenía y el costo era alto, porque una entrada nueva costaba dos
publicaciones y una de ellas era una versión de la herramienta cuyo único
cambio sustantivo era un número. El determinismo que RF-05 pide no lo daba la
versión del catálogo sino el pin de la plantilla, porque lo que determina la
estructura de un producto es el contenido del activo del que se deriva, no la
lista de activos que existen. El catálogo pasa entonces a resolverse en
ejecución, con tres fuentes en orden —la ruta que dé la invocación, la última
versión que sirve el registro, y la copia instalada cuando el registro no
responde—, informando cuál de ellas respondió.

Lo que sí queda fijado dentro de la herramienta es la versión exacta de cada
plantilla que esa entrega ofrece, con una versión exacta y nunca un rango,
porque un rango haría que dos ejecuciones de la misma herramienta resolvieran
versiones distintas. Un módulo no se fija así, sino que se resuelve a la
última publicada y la versión obtenida queda registrada en el producto, ya
que pinar cada módulo obligaría a publicar una versión de la herramienta por
cada versión de cada módulo, que es exactamente la dependencia que el
catálogo existe para evitar.

\section{El núcleo de la herramienta y la operación de derivación}
\label{sec:herramienta-cli}

Con los activos construidos y el catálogo en operación, el hito siguiente es
el mecanismo que los consume. La herramienta se construyó separando lo que
toda operación comparte de lo que es propio de una, y esa separación es lo
que permitió a la fase siguiente ampliarla en lugar de rehacerla.

\subsection{El núcleo y sus puertos}
\label{sec:nucleo-herramienta}

Pertenece al núcleo aquello que una segunda operación necesitaría igual: la
resolución de la entrada contra el modelo de variabilidad, la consulta al
catálogo, la lectura y escritura del registro de versiones del proyecto, la
conversión de toda escritura en un plan de cambios antes de aplicarlo, y la
traducción de cualquier fallo a un código de salida estable.

El marco exige arquitectura hexagonal al proyecto generado, y la herramienta
la aplica a sí misma. Su núcleo no conoce el sistema de archivos, el registro
de paquetes, la terminal ni el reloj, sino que los declara como los seis
puertos de la Tabla~\ref{tab:puertos-nucleo} y recibe sus implementaciones
desde fuera.

\begin{table}[H]
\centering
\begin{tabular}{p{4.8cm}p{10.2cm}}
\toprule
\textbf{Puerto} & \textbf{Qué abstrae} \\
\midrule
Sistema de archivos & Leer, escribir y comprobar existencia de rutas \\
Registro de paquetes & Obtener un activo y resolver la dirección del registro de un ámbito \\
Preguntas & Pedir al desarrollador lo que la invocación no trajo \\
Reporte & Lo que la herramienta dice, con su forma \\
Reloj & El momento que el producto registra como su procedencia \\
Ejecutor de comandos & Invocar procesos externos, como el gestor de paquetes o el control de versiones \\
\bottomrule
\end{tabular}
\caption{Los seis puertos del núcleo de la herramienta}
\label{tab:puertos-nucleo}
\end{table}

La consecuencia práctica es que el comportamiento del núcleo se prueba sin
tocar disco, sin red y sin terminal, con dobles escritos a mano que
implementan cada puerto. Esa decisión es la que hace que la suite corra en
segundos y que una prueba falle por el motivo que dice y no por el entorno
donde se ejecuta.

Ese contrato es estrecho, porque una operación declara su nombre, su resumen,
los parámetros que admite y una función que recibe la entrada ya validada
junto con el contexto de puertos.
El despachador no sabe qué operaciones existen, sino que las lee de un
registro, de modo que agregar una es escribir su módulo y registrarla sin
modificar el núcleo. Ese punto de extensión es el entregable que la fase
siguiente consume, y su existencia se comprobó antes de necesitarla, porque la
operación de consulta del catálogo se construyó en esta misma fase sobre el
mismo contrato, de modo que la segunda operación existió mientras el punto de
extensión todavía podía corregirse.

\subsection{La operación de derivación}
\label{sec:comando-generacion}

La operación de derivación resuelve la plantilla contra el catálogo y la
obtiene del registro, recoge las decisiones de variabilidad, comprueba que la
combinación elegida satisfaga las reglas del modelo, y convierte todo ello en
un plan de cambios que el núcleo escribe. Al terminar existe un proyecto que
arranca sin que nadie edite un archivo.

Una operación de negocio elegida no produce un archivo sino un recorrido
completo desde la superficie externa hasta el contrato del repositorio: un
método en el adaptador HTTP con su ruta, un caso de uso en la capa de
aplicación, el método que ese caso de uso necesita en el puerto del
repositorio, y un evento de dominio cuando la operación cambia el estado de la
entidad que la capacidad gobierna. Ese reparto no se decide archivo por
archivo, sino que cada operación recorre el mismo camino desde que se elige
hasta los archivos que produce, pasando por la entrada que la herramienta
resuelve y por las reglas que convierten un término de acción en un nombre de
clase, una ruta, un método y un evento. Es lo que hace que la selección de
características tenga consecuencias estructurales y no cosméticas.

\paragraph*{El límite de lo que se genera}

El mecanismo produce la estructura, los nombres y los contratos vacíos,
mientras que las reglas de negocio las construye el equipo que recibe el
proyecto, y mantener esa distinción es lo que mantiene al artefacto dentro de
su alcance.

La distinción se declara donde se nota, de modo que las operaciones cuyo
efecto se sigue del propio término de acción se generan completas; las que
dependen de la capacidad se generan cargando la entidad, persistiéndola,
publicando su evento y dejando el efecto marcado con un comentario. El
registro del producto declara, por operación, cuál de los dos casos es, y el
producto lo dice en su propia documentación. Un producto que pareciera más
terminado de lo que está sería el peor resultado posible de un generador.

\paragraph*{Determinismo}

La misma entrada produce el mismo producto, sin importar quién ejecute el
comando. Tres decisiones lo sostienen: la versión de la plantilla se resuelve
contra el pin que viaja en la herramienta y no contra el registro; toda
escritura pasa por un plan de cambios que es un dato inspeccionable antes de
aplicarse; y el único valor que varía entre dos ejecuciones idénticas, que es
el momento de la derivación, se obtiene por un puerto y se registra en un
campo conocido. La verificación deriva el mismo producto dos veces y compara
los árboles resultantes excluyendo ese campo.

\paragraph*{El recorrido de entrada}

La operación cierra con una revisión que muestra cada decisión tomada y
ofrece corregir cualquiera de ellas antes de escribir nada. Corregir una
respuesta vuelve a preguntar únicamente lo que depende de ella, y la razón es de
adopción y no de comodidad, porque una pregunta del estándar contestada mal
costaba el comando entero, y ese costo es el que hace que un desarrollador
deje de leer las preguntas. Fuera de una terminal la revisión declara y
continúa, porque una corrida automatizada no tiene a quién preguntar.

El recorrido completo se corrigió además en un punto que ninguna prueba veía,
y la causa es estructural, porque los circuitos automáticos corren con
banderas y ningún camino automático pasa por una pregunta. Lo que apareció al
recorrerlo a mano fue que una condición no satisfecha se comprobaba después
de cerrar todas las preguntas, y el costo de no satisfacerla era perderlas
todas. Hoy la condición viaja con la pregunta y se rechaza ahí mismo, la
elección de plantilla ocurre antes de la primera pregunta del modelo, y la
tecla de escape retrocede una pregunta en lugar de abandonar el comando.

El recorrido abre eligiendo la plantilla, antes de cualquier pregunta del
modelo, ofreciendo las dos opciones del almacén con lo que cada una deriva
como pista y mostrando la lista entera en lugar de enrollarse sobre sí misma.

\section{Verificación y conformidad}
\label{sec:pruebas-fase-2}

El último hito no agrega funcionalidad, sino que produce la evidencia de que
lo construido hace lo que debía, y opera en dos niveles cuya distinción es
deliberada. El primero prueba la herramienta con dobles de sus puertos,
comprobando la resolución de la entrada, la evaluación de las reglas del
modelo, el reparto de archivos por operación y los nombres derivados, sin
tocar disco ni red. El segundo deriva productos reales y los somete a su
propia verificación —análisis estático, pruebas y compilación—, que es la
única forma de saber que lo generado es código válido y no texto bien
formado.

\subsection{Romper la regla para saber que existe}

Un producto que supera su análisis estático demuestra que el analizador
corre, no que la regla de dependencia esté activa. La verificación introduce
por eso, a propósito, una importación que la regla prohíbe, exige que el
análisis falle nombrando esa regla, y restaura el archivo. El criterio se
generalizó a toda regla nueva de la línea: una regla cuya violación no se
haya comprobado no está verificada, está escrita.

El criterio dio su propio resultado, porque una regla que exigía una prueba de
humo por cada ruta declarada se satisfacía con que la ruta apareciera en el
título de la prueba, no con que alguna petición la ejercitara, de modo que una
prueba apuntada a otra ruta la habría pasado. El defecto apareció al intentar
romperla y no al escribirla.

\subsection{Conformidad del producto derivado}

El producto se verifica contra la condición de llegada de la fase y no contra
una lista de archivos: arranca sin edición previa, responde sus sondas,
registra lo que sirve, no registra las sondas, y no contiene direcciones ni
lecturas del entorno fuera de su módulo de configuración. El circuito deriva
varios productos con entradas distintas —capacidades distintas, las dos
formas de ruta, con y sin eventos, una capacidad declarada fuera del
catálogo— y cada uno pasa su propia verificación, que es la evidencia que el
Capítulo~\ref{cap:evaluacion} cita.

La documentación de instalación y uso de la herramienta vive en el
repositorio de los activos y viaja en el paquete publicado, no en este
documento, por una razón de mantenimiento, y es que una guía de uso transcrita
en la tesis quedaría desactualizada en la primera versión publicada.

\section{Resumen del capítulo}

Esta fase convirtió un encuadre en un mecanismo, de modo que lo que al cerrar el
Capítulo~\ref{cap:fundamentacion} era una definición es ahora una herramienta
instalable que produce proyectos conformes a ella sin intervención manual, y
con ella las tres reglas concretas que el marco había dejado asignadas: la
forma de un producto válido, los contratos por los que se le acopla una
capacidad compartida, y el gobierno de los activos que lo sostiene en el
tiempo.

Cuatro decisiones estructurales sostienen el resultado, y la primera es el
orden de los hitos, porque los activos se construyeron antes que el
generador, de modo que la plantilla base funcionó como validación ejecutable
del marco cuando corregirlo todavía era barato. La segunda es la separación
entre el núcleo común y la operación concreta, con un punto de extensión
comprobado antes de necesitarlo, que es lo que permite a la fase siguiente
ampliar la herramienta en lugar de rehacerla. La tercera es que la plantilla
base es la salida de sus propias plantillas y no un artefacto mantenido en
paralelo, lo que hace imposible que se separen en silencio. Y la cuarta es el
\emph{contract package} como artefacto neutral del que dependen productos y
módulos sin depender entre sí, que es la pieza de la que depende todo el
mecanismo de composición del capítulo siguiente.

La variabilidad quedó declarada en artefactos versionados y no en el código
del generador, y esa separación es la que permitió después ampliar el espacio
de selección, incorporar el segundo nivel del estándar y admitir una
capacidad fuera del catálogo sin modificar la herramienta. Es también el
punto donde el estándar interviene en el mecanismo y no solo en el registro,
porque sin el espacio de selección la alineación con BIAN se reduciría a un
metadato escrito a mano.

El quinto hito dejó la lección más incómoda y la más útil, porque con un solo
activo base, el reparto entre lo que la plantilla declara y lo que la
herramienta sabe podía no decidirse, y la consecuencia fue que la herramienta
guardaba el catálogo de archivos del producto dentro de su código. Cerrar la
promesa de un almacén obligó a hacer ese reparto explícito, y lo que salió es
más firme que lo que había: una plantilla que se describe a sí misma, un motor
cuya capacidad de decisión está enumerada, una derivación que es un intérprete
y una forma de producto comprobable por los dos lados. La evidencia de que el
resultado es genérico no es un argumento sino una segunda plantilla, de otro
modelo de ejecución, que entró sin que una línea de la herramienta cambiara.

Queda abierto, por construcción, todo lo que exige módulos que aún no
existen: incorporarlos a un producto, consultarlos, actualizarlos y preservar
el código propio del equipo al hacerlo. Es el objeto del
Capítulo~\ref{cap:gestor}.


%%----------------------------------------------------------
%%  CAPÍTULO 5: DISEÑO Y CONSTRUCCIÓN DEL GESTOR DE MÓDULOS REUTILIZABLES
%%----------------------------------------------------------

\chapter{Diseño y construcción del gestor de módulos reutilizables}
%%==========================================================
%%  CAPÍTULO 5: DISEÑO Y CONSTRUCCIÓN DEL GESTOR DE MÓDULOS REUTILIZABLES
%%==========================================================
\label{cap:gestor}

El Capítulo~\ref{cap:generacion} dejó una herramienta que produce un proyecto
conforme al marco y un contrato por el que una capacidad compartida puede
acoplársele. Lo que no existía todavía era esa capacidad, ni la operación que
la incorpora. Un proyecto generado nace completo el primer día y empieza a
quedarse corto el segundo, cuando necesita algo que no eligió al derivarse, y
esta fase construyó lo que resuelve ese segundo día.

Lo que entregó son dos cosas que se validan mutuamente, y la primera son
módulos reales del catálogo, que es la única forma de comprobar que el
contrato de integración sirve para algo. La segunda son las operaciones que
los incorporan, los consultan, los actualizan y los retiran sobre un proyecto
que ya existe y que su equipo ya modificó. La condición de llegada se formuló
antes de construir, y exige que un producto derivado al que se le incorpora un
módulo siga arrancando sin edición manual, supere su propia verificación,
conserve el código que su equipo escribió y declare qué módulo tiene y en qué
versión. El detalle de la construcción se conserva en los documentos de la
Fase 3, disponibles en la carpeta \texttt{Implementación/} del repositorio del
trabajo de grado.

\section{Hitos de la construcción}
\label{sec:hitos-fase-3}

La fase se ejecutó en cuatro hitos, y el orden entre ellos importa por la
misma razón que en la fase anterior, porque un módulo real contra la
definición encuentra lo que ninguna lectura de la definición encuentra. Por
eso el activo componible y su prueba contra un caso real preceden a las
operaciones que lo manejan, y la evidencia del circuito completo cierra la
fase. La Tabla~\ref{tab:hitos-gestor} recoge qué dejó cada hito y cómo quedó
comprobado.

\begin{table}[H]
\centering
{\footnotesize
\begin{tabular}{p{3.1cm}p{5.9cm}p{5.4cm}}
\toprule
\textbf{Hito} & \textbf{Qué dejó construido} & \textbf{Cómo quedó comprobado} \\
\midrule
Qué es un módulo & El manifiesto con que un módulo se describe a sí mismo, la frontera entre lo empaquetado y lo generado, y las condiciones que lo hacen válido & La herramienta incorpora un módulo sin conocer su nombre, leyendo solo lo que el manifiesto declara \\
El primer módulo & La persistencia relacional, con su puerto en el contrato y su adaptador en el producto & El producto compuesto guarda un dato, se reinicia y lo vuelve a leer \\
Las operaciones & Incorporar, consultar, actualizar y retirar, sobre el punto de extensión de la fase anterior & El núcleo creció sin que ninguna operación anterior cambiara \\
El segundo módulo & La observabilidad, de otra naturaleza que la primera & Entró sin cambiar el núcleo, el esquema del manifiesto ni la operación \\
\bottomrule
\end{tabular}
}
\caption{Hitos de la construcción del gestor de módulos}
\label{tab:hitos-gestor}
\end{table}

La frontera con la fase siguiente es la que separa construir de medir. Esta
fase comprueba que el mecanismo funciona; si funciona en manos de un equipo
que no lo construyó, y cuánto esfuerzo ahorra, es lo que mide el
Capítulo~\ref{cap:evaluacion}.

\section{Las capacidades transversales y la que se construyó}
\label{sec:capacidades-transversales}

El inventario de capacidades partió del análisis cualitativo del contexto y
no de una taxonomía externa. Aquí hay una desviación respecto del
anteproyecto que debe declararse, y es que aquel documento pedía categorizar
las capacidades transversales tomando como referencia los dominios funcionales
del estándar de negocio, y no se hizo, porque ese estándar modela capacidades
bancarias y no infraestructura técnica de aplicaciones. Categorizar un
registro de eventos contra un dominio de negocio habría producido una
clasificación sin poder explicativo.

De las capacidades inventariadas se eligió una, con criterios escritos antes
de aplicarlos: que fuera verificable por comportamiento observable, que no
dependiera de infraestructura externa al entorno de verificación, que
ejercitara el contrato de composición en su forma completa, que su ausencia
fuera visible en el producto, y que no exigiera decisiones de seguridad que
el contexto no hubiera fijado.

La elegida fue la persistencia relacional del producto, y la
elección corrigió una afirmación anterior del propio proyecto, que daba por
decidido que el módulo comprometido sería el de cifrado y manejo de tokens.
El criterio descarta el cifrado precisamente por el último punto, ya que
exige fijar antes decisiones de seguridad, como la gestión de claves, su
rotación y su custodia, que pertenecen a la organización y no a una línea de
producto.

El inventario fijó además un orden de crecimiento, y la primera capacidad de
ese orden se construyó después como segundo módulo
(sección~\ref{sec:segundo-modulo}). Esa segunda construcción es la que
permite afirmar algo que con un solo módulo no era afirmable.

\section{Qué es un módulo}
\label{sec:anatomia-modulo}

Un módulo de esta línea es un paquete publicado que implementa un puerto
técnico del \emph{contract package} y que se describe a sí mismo. La
herramienta no conoce ningún módulo por su nombre, sino que lee un manifiesto
que viaja dentro del paquete y que declara qué puerto provee, qué
configuración necesita, qué archivos aporta al producto y en qué puntos de
este hay que insertar su cableado.

La frontera entre lo que viaja en el paquete y lo que se escribe dentro del
producto tiene una consecuencia que vale más que la elegancia, porque
incorporar un módulo nuevo al catálogo no exige publicar una versión nueva de
la herramienta. Sin ella, cada módulo habría requerido código propio en el
generador, y el catálogo habría quedado acoplado a su versión.

\subsection{La frontera entre lo empaquetado y lo generado}

La distinción que gobierna todo el mecanismo es cuál de las dos partes de un
módulo vive dónde. La \emph{implementación} viaja empaquetada y versionada en
el módulo, y es código que el equipo del producto no edita y que una
actualización reemplaza por completo. El \emph{cableado} se genera dentro del
producto, y comprende el adaptador que ata el puerto a la implementación, la
configuración que el producto valida junto con la suya, y el registro de qué
quedó incorporado.

De ahí se sigue la regla más restrictiva de la fase: un módulo aporta
plantillas, nunca clases base. Si el módulo exportara una clase de
la que el producto heredara, el código del producto quedaría atado a la forma
interna del módulo y una actualización podría romperlo sin tocarlo. Lo
generado es del equipo que lo recibe, que es lo que RNF-02 exige.

Para que el cableado pueda entrar, y sobre todo para que pueda reemplazarse
después, el producto trae puntos de inserción declarados: unos son marcas de
una sola línea donde se agrega algo, y otros son regiones delimitadas cuyo
contenido se sustituye entero. La diferencia importa, porque lo que se
inserta sin delimitar no se puede reemplazar, y eso apareció al construir la
actualización y no al diseñar la composición.

Un módulo escribe también en la documentación del producto, en una región que
él declara. La razón es que el producto dice de sí mismo cosas que dejan de
ser ciertas al componerlo. Un producto que advierte que pierde sus datos al
detenerse deja de perderlos cuando entra el módulo de persistencia, y esa
frase tiene que corregirse por construcción y no por memoria de quien compone.

\subsection{A qué forma de producto sirve}

El manifiesto declara además a qué forma de producto sirve el módulo, y
mientras la línea tuvo una sola plantilla, esa declaración habría sido
redundante; cuando el almacén admitió una segunda forma dejó de serlo, porque
el cableado de un módulo aterriza en puntos de inserción concretos y una
forma que no los tenga lo recibiría roto.

La comparación se hace contra el registro del producto y nunca contra la
plantilla que lo derivó, porque esa plantilla puede no estar instalada, y se
hace por identificador y versión entera y nunca por rango, ya que una versión
de forma se mueve exactamente cuando un punto se mueve o desaparece. Ser
indulgente ahí convertiría un rechazo que nombra la causa en un producto roto
sin aviso.

Las dos declaraciones que deciden una incorporación caben en cinco líneas del
manifiesto, y son el puerto que el módulo provee y la forma a la que sirve.
Todo lo que la herramienta necesita saber de un módulo para aceptarlo o
rechazarlo está ahí, y nada de ello está en la herramienta.

El mensaje del rechazo nombra las dos formas, la que el módulo sirve y la que
el producto tiene, y explica que la incompatibilidad no es cuestión de
esfuerzo sino de que el cableado no tendría dónde aterrizar. El producto queda
sin tocar.

\section{El módulo de persistencia como caso de referencia}
\label{sec:modulo-persistencia}

El módulo se construyó sobre el motor embebido que la plataforma de ejecución
ya trae, sin dependencias de ejecución nuevas. Esa restricción no es
estética, porque una dependencia de ejecución en un entorno regulado abre una
auditoría, y el objetivo era probar el mecanismo de composición y no la
elección de un motor.

La persistencia es además el caso que obliga a ejercitar a la vez las dos
clases de puerto que el Capítulo~\ref{cap:generacion} definió. El repositorio
del agregado es un puerto de dominio, que el producto declara y realiza. El
acceso al motor es un puerto técnico, que el \emph{contract package} declara y
el módulo realiza. El adaptador del repositorio, que vive en el producto,
consume el puerto técnico del módulo, y ninguna dependencia une directamente
al producto con el módulo.

El módulo aporta la capacidad técnica pero no una base de adaptador que el
producto extienda, porque la herencia entre un activo del catálogo y el
código del producto crearía un acoplamiento que ninguna de las dos partes
controla. La consecuencia es que el dominio no conoce el motor de
persistencia, y cambiarlo no altera ninguna declaración de la capa de
dominio. De aquí se desprende una regla adicional de validez, según la cual
un producto no puede tener puertos de dominio sin adaptador, y es la que
después impide retirar un módulo que dejaría al producto incompleto.

\paragraph*{Un hallazgo del contenedor de inyección}

Un módulo no alcanzaba el puerto de configuración del producto, porque el
contenedor de dependencias resuelve dentro del módulo que pide. La solución
fue que el producto exponga los puertos de la línea en un módulo global
propio, separado de su raíz de composición. Es una consecuencia del marco que
ninguna lectura del marco habría anticipado, porque apareció al ejecutar, y es
el tipo de hallazgo que justifica construir un módulo real en lugar de
especificar cómo sería.

\section{Las cuatro operaciones del gestor}
\label{sec:gestor-modulos}

Las operaciones del gestor se montan sobre el punto de extensión que la fase
anterior dejó, y ninguna lo modifica. El núcleo sí creció, y la distinción
entre crecer y modificarse es la que mantuvo la frontera entre las dos fases.

\subsection{Incorporar un módulo a un proyecto existente}
\label{sec:composicion}

La composición resuelve el módulo contra el catálogo, lo obtiene del
registro, lee su manifiesto y convierte lo que el manifiesto declara en un
plan de cambios sobre un producto que ya existe. El plan sigue siendo el
mismo dato inspeccionable de la fase anterior; lo que cambió es que ahora
puede escribir sobre un destino que no está vacío.

Escribir sobre un proyecto ajeno exigió tres cosas que generar no
necesitaba: un modo que admite un destino existente, dos clases de cambio
—reescribir un archivo completo y reemplazar el contenido de una región
delimitada— y respaldo de todo lo que se toca antes de tocarlo. El diseño
original del plan de cambios ya había anticipado la necesidad, lo que
permitió extenderlo en lugar de reemplazarlo.

\paragraph*{Lo que la operación se niega a hacer}

Antes de producir un solo cambio, la operación comprueba nueve reglas sobre
el producto y sobre el módulo. Rechaza cuando el módulo no sirve a la forma
del producto, cuando el destino no es un producto derivado, cuando el activo
no está en el catálogo o no está activo, cuando el módulo ya está
incorporado, cuando otro módulo ya provee el mismo puerto, cuando las
versiones de contrato no coinciden, cuando falta un punto de inserción que el
cableado necesita, y cuando un archivo que el módulo genera ya existe en el
producto.

Que las nueve se comprueben antes y no durante la escritura es lo que hace
que un rechazo deje el producto exactamente como estaba, y esa propiedad se
verifica. Un rechazo que dejara el producto a medio modificar sería peor que
no tener la operación. La razón de rechazar en lugar de intentar arreglar es
que cada regla describe un producto que la herramienta no sabe componer sin
suponer algo sobre código que no escribió, y suponer en código ajeno es el
modo más caro de equivocarse, porque el daño aparece después, en un producto
que ya nadie relaciona con el comando que lo tocó.

\paragraph*{Lo que apareció al ejecutar}

Los tres defectos de la Tabla~\ref{tab:defectos-composicion} no aparecieron
al diseñar la composición sino al correrla, y son la parte más instructiva de
la fase.

\begin{table}[H]
\centering
{\footnotesize
\begin{tabular}{p{4.9cm}p{9.8cm}}
\toprule
\textbf{Defecto} & \textbf{Qué lo causaba} \\
\midrule
Un archivo tocado dos veces perdía el primer cambio & El plan aplicaba cada cambio sobre el archivo original y no sobre el resultado del anterior \\
Un marcador se confundía con su propia mención en la documentación del código & La búsqueda del punto de inserción no distinguía la marca de un comentario que la nombra \\
El producto compuesto compilaba pero no pasaba su análisis estático & Una importación quedaba sin uso al reemplazar una región \\
\bottomrule
\end{tabular}
}
\caption{Defectos detectados al ejecutar la composición}
\label{tab:defectos-composicion}
\end{table}

El tercero es el que mejor justifica el circuito de verificación completo,
porque un producto compuesto que compila pero no pasa su análisis estático es un
producto que su equipo no puede entregar, y ninguna prueba unitaria de la
herramienta lo habría mostrado.

\subsection{Consultar, actualizar y retirar}
\label{sec:consulta-actualizacion}

La consulta responde qué ofrece el catálogo, qué tiene incorporado el
producto, en qué versión, cuál quedó por detrás de la disponible y qué archivo
generado cambió el equipo. Es lo que permite decidir si actualizar antes de
intentarlo, y dice además lo que el producto \emph{tiene} y no lo que podría
tener, que es una corrección posterior. Antes declaraba disponible un módulo
que la composición habría rechazado, porque no comprobaba la forma del
producto.

La actualización mueve un módulo incorporado a una versión posterior sin
intervención manual, y su problema real es qué hacer con lo que el equipo
cambió. La respuesta fue comparar una huella del archivo generado contra la
que el registro guardó: si coinciden, el archivo se reemplaza; si no, se
preserva y la herramienta lo dice. No se fusiona nada, porque fusionar
exigiría que la herramienta decidiera sobre código que no escribió, y la
decisión de informar en lugar de resolver es lo que mantiene el control en el
equipo. La actualización comprueba además tres de las nueve reglas de la
composición, que al principio no repetía, de modo que un cambio de versión
podía mover un módulo a una versión construida contra otra mayor del contrato,
o contra otra forma de producto, sin que nadie lo mirara.

A las tres operaciones se sumó después una cuarta, y es la única pieza del
artefacto que ningún requerimiento pide. Se construyó por decisión
de alcance del autor, declarada como tal, con el argumento de que una
decisión que no se puede deshacer no es una decisión, porque un equipo que no
puede retirar un módulo no lo incorpora para probarlo. La operación lee
únicamente el registro del producto, restaura la región que el módulo había
reemplazado y conserva lo que el equipo editó; y rechaza antes de tocar nada
si al retirarlo un puerto de dominio quedaría sin adaptador. Queda un límite
declarado, porque un archivo generado que el equipo modificó se conserva por
la misma razón que lo conserva la actualización, y ese archivo sigue
importando el módulo que acaba de salir. La herramienta lo dice en lugar de
corregirlo.

\section{El segundo módulo del catálogo}
\label{sec:segundo-modulo}

Un módulo construido demuestra que el mecanismo funciona para ese módulo, pero
no que sea general, porque con un solo caso no hay forma de distinguir un
contrato de composición de una solución escrita para él. El segundo módulo es
la evidencia de esa diferencia, y por eso se eligió de otra naturaleza que el
primero.

La capacidad construida fue la observabilidad, que era la primera
del orden de crecimiento declarado en el inventario. Donde el primer módulo
persiste datos, este los \emph{exporta}, entregando a un colector externo las
entradas del registro, el tramo de traza de cada petición y su duración, en
un protocolo abierto de telemetría y sin dependencias de ejecución nuevas.
Quedó fuera de su alcance exponer un endpoint para que un sistema externo
recolecte métricas, y la razón es del marco y no del módulo, porque un módulo
aporta adaptadores de salida y un endpoint es superficie de entrada.

\paragraph*{Qué quedó probado}

Un módulo de otra naturaleza entró sin cambiar el núcleo de la herramienta, ni
el esquema del manifiesto, ni la operación de composición. Lo único que creció
fue el catálogo, con una entrada, que es la afirmación que el primer módulo no
podía sostener.

El contrato ganó además un puerto con una propiedad que ninguno anterior
tenía, porque es el único que ningún adaptador de la plantilla
satisface. El producto reporta sus señales exactamente igual tenga o no el
módulo incorporado, y lo inyecta de forma opcional, de modo que incorporar el
módulo no cambia una línea del producto. La plantilla tampoco reemplaza su
registro de eventos cuando el módulo está presente, porque el asiento es del
producto, sino que entrega al módulo la misma entrada que ya escribió.

\paragraph*{Dos hallazgos que solo aparecen al ejecutar}

El primero es una regla del marco que nadie había enunciado, y es que un
módulo no puede recibir el puerto al que la plantilla le reporta. El adaptador
pedía el puerto de registro de eventos para escribir sus propias trazas, lo
que cerró un ciclo con el registro del producto; el resultado fue que el
producto arrancó, atendió peticiones y \emph{no exportó nada, en silencio},
que es la peor forma de fallar. La corrección dejó una prueba de regresión en
la plantilla, porque una regla cuya violación no se comprueba no está
verificada sino escrita.

El segundo corrigió una afirmación del propio marco, que se sostenía que una
versión menor del contrato llega a los consumidores sin intervención, y en el
tramo de versiones anterior a la primera estable eso es falso, porque el
rango que un paquete declara no acepta la menor siguiente. La consecuencia
fue que el primer módulo tuvo que republicarse sin un solo cambio de código.
Es exactamente el costo que el marco había anticipado al versionar el
contrato como una unidad, apareciendo por primera vez en ejecución.

\section{Verificación y conformidad}
\label{sec:verificacion-composicion}

La verificación de esta fase no podía ser una suite de pruebas unitarias. Lo
que había que comprobar es un circuito que atraviesa varias herramientas y
dos procesos: derivar un producto, incorporarle el módulo, verificar que el
producto sigue pasando su propia verificación, arrancarlo, almacenar un dato,
detener el proceso, volver a arrancarlo y leer el dato de vuelta. Eso se
construyó como un circuito automatizado que corre en la puerta de
verificación de la línea.

El circuito comprueba además cada rechazo, dejando el producto idéntico en
todos ellos, la actualización preservando el archivo que el equipo modificó,
y el producto volviendo a ser conforme cuando se le retira el último módulo.
Incorpora el mismo módulo a dos productos distintos, con agregados distintos,
y comprueba que el cableado se nombra por el agregado de cada uno, que es la
evidencia de que un activo se comparte entre proyectos sin acoplarlos. Y
porque el segundo producto se compone después de actualizar el primero, los
dos quedan con versiones distintas del mismo módulo a la vez, que es el
criterio literal de RNF-06.

La evidencia se cita en los documentos de la Fase 3 con su fecha y con el
comando que la reproduce; este capítulo no la reproduce, porque una
transcripción de salidas en la tesis queda desactualizada en la primera
versión publicada.

El criterio de la fase anterior se aplicó a todo lo que esta agregó, y dice
que una regla cuya violación no se haya comprobado no está verificada. Cada regla
nueva se rompió a propósito, se comprobó que el circuito fallara con su
propio mensaje, y se restauró.

\section{Resumen del capítulo}

Esta fase cerró el artefacto, de modo que al terminar existe una herramienta
que genera un proyecto conforme al marco, le incorpora una capacidad
transversal sobre un proyecto ya en curso, informa de su estado, la actualiza
preservando lo que su equipo escribió y la retira si hace falta.

Tres decisiones explican el resultado, y la primera es que el módulo se
describe a sí mismo, lo que desacopla el catálogo de la versión de la
herramienta y hace que incorporar un módulo nuevo no sea un cambio de código.
La segunda es la frontera entre la implementación empaquetada y el cableado
generado, con la regla de que un módulo aporta plantillas y nunca clases
base, que es lo que hace que lo generado sea del equipo y que una
actualización no pueda romper código que no reemplaza. La tercera es que la
actualización informa en lugar de fusionar, lo que mantiene en el equipo toda
decisión sobre código que la herramienta no escribió.

La fase dejó además una lección metodológica que ninguna lectura del marco
habría producido, porque varias de sus decisiones corrigen defectos que
aparecieron al \emph{ejecutar} y no al diseñar. El más incómodo y el más útil
lo dio el segundo módulo, donde un adaptador pedía el puerto de registro de
eventos al que la propia plantilla le reporta, cerró un ciclo, y el producto
arrancó,
atendió peticiones y no exportó nada en silencio. Un fallo silencioso es el
que ninguna lectura del marco produce y ninguna prueba del módulo aislado
detecta, y solo apareció al componer el módulo en un producto real y mirar
qué salía por el otro lado.

Ese segundo módulo es también lo que convierte una afirmación en evidencia, porque
con un módulo construido, el mecanismo funciona para ese módulo; con uno de
otra naturaleza, que exporta en lugar de persistir y que entró sin cambiar el
núcleo, ni el esquema del manifiesto, ni la operación de composición, lo que
queda probado es el mecanismo. Es la misma clase de evidencia que la segunda
plantilla aporta del lado de la derivación, y las dos juntas son lo que
permite hablar de una línea de productos y no de un generador con un ejemplo.

Lo que queda no es construcción sino medición, es decir cuánto del esfuerzo de
configuración inicial se reduce realmente, si la composición funciona en
manos de un equipo que no construyó la herramienta, y si un proyecto derivado
permite empezar a escribir lógica de negocio antes que el proceso actual. Es
el objeto del Capítulo~\ref{cap:evaluacion}.


%%----------------------------------------------------------
%%  CAPÍTULO 6: EVALUACIÓN Y VALIDACIÓN DEL ARTEFACTO
%%----------------------------------------------------------

\chapter{Evaluación y validación del artefacto}
%%==========================================================
%%  CAPÍTULO 6: EVALUACIÓN Y VALIDACIÓN DEL ARTEFACTO
%%==========================================================
\label{cap:evaluacion}

\section{Diseño de la evaluación}

\section{Resultados de la evaluación}

\section{Resumen del capítulo}

%%----------------------------------------------------------
%%  CAPÍTULO 7: CONCLUSIONES
%%----------------------------------------------------------

\chapter{Conclusiones}
%%==========================================================
%%  CAPÍTULO 7: CONCLUSIONES
%%==========================================================

Las conclusiones deben ser la respuesta a los objetivos o propositos planteados así como presentar las lecciones aprendidas más importantes durante el proceso. Este capitulo debería incluir: 

\begin{itemize}
\item \textbf{Resumen de los hallazgos}: Las conclusiones deben proporcionar un resumen claro y conciso de los principales hallazgos de la investigación. Esto debería incluir una revisión de los objetivos de la investigación y una discusión de cómo se cumplieron estos objetivos.
\item \textbf{Relación con la literatura existente}: Las conclusiones deben relacionar los hallazgos de la investigación con la literatura existente. Esto podría implicar la comparación de los resultados con estudios previos, la discusión de cómo los hallazgos contribuyen a la teoría existente, etc.
\item \textbf{Implicaciones prácticas}: Si es relevante, las conclusiones deben discutir las implicaciones prácticas de los hallazgos de la investigación. Esto podría incluir recomendaciones para la práctica profesional, sugerencias para la implementación de políticas, etc.
\item \textbf{Limitaciones de la investigación}: Las conclusiones deben reconocer las limitaciones de la investigación. Esto podría incluir limitaciones en la metodología, limitaciones en la recopilación de datos, etc.
\item \textbf{Sugerencias para futuras investigaciones}: Las conclusiones deben proporcionar sugerencias para futuras investigaciones. Esto podría incluir áreas que no se exploraron en la investigación actual, preguntas que surgieron como resultado de la investigación, etc.
\item \textbf{Claridad y coherencia}: Las conclusiones deben estar bien escritas y presentadas de manera clara y coherente. Deben proporcionar un cierre efectivo para la tesis.
\end{itemize}

\section{Conclusiones}

\section{Trabajos futuros}

\section{Lecciones aprendidas}

\section{Puede agregar más secciones si lo necesita}

%%----------------------------------------------------------
%%  BIBLIOGRAFÍA
%%----------------------------------------------------------

\bibliographystyle{apalike}
\bibliography{biblio}

\end{document}